\partdivider{C}{Thermodynamics and the Molecular Picture}{The laws governing heat, work and energy, and the kinetic theory that explains them from the ceaseless motion of molecules. Here the whole volume converges.}{\chref{ch:thermo} Thermodynamics\\[1.2mm]
\chref{ch:kinetic} Kinetic Theory}
\ncertchapter{4}{Thermodynamics}\chlabel{ch:thermo}

\chref{ch:thermal} described what heat \emph{does} to matter: it raises temperatures, expands solids,
melts ice and flows from hot to cold. This chapter asks a sharper question. When heat and work are
exchanged, what rules govern the conversion of one into the other? A kilogram of coal releases about
\SI{30}{MJ} of heat when it burns; how much of that can a power station turn into electrical work, and
why can it never be all of it?

The answers are the laws of thermodynamics, and they are unusual among physical laws in needing no
model of what matter is made of. The \textbf{first law} is energy conservation, extended to count heat as
a transfer of energy: it is the bookkeeping of \chref{ch:energy}, where friction's ``lost'' mechanical
energy reappeared as heat, now run in both directions. The \textbf{second law} is new. It says which
energy-conserving processes actually happen, and it sets a ceiling on the efficiency of every engine that
depends only on two temperatures. From it comes \textbf{entropy}, the quantity that tells the future from
the past.

The method throughout is the one used for the gas laws: describe a system by a few measurable state
variables (pressure, volume, temperature), and follow how they change. We use the ideal gas as the
working example because its equation of state $PV=nRT$ (\chref{ch:thermal}, Section~\ref{sec:gaslaw}) makes every integral
doable; where a result holds only for the ideal gas, we say so. \chref{ch:kinetic} then opens the box
and explains the same results from molecules in motion.

\begin{motivbox}{Why thermodynamics matters}
Thermodynamics is the physics of engines and of limits; it built the industrial world and drew the
boundary of what is possible.
\begin{itemize}
\item \textbf{The first law is energy accounting for heat and work.} It tells you how much fuel an engine
needs, how a refrigerator moves heat ``uphill'' at a cost, and why you cannot get more work out than
energy in. Every power station, engine and air conditioner is a first-law balance sheet.
\item \textbf{The second law sets a ceiling no cleverness can beat.} Carnot's result, that even a perfect
engine must reject a fixed fraction of its heat, set only by the two temperatures, is why car engines
manage 25 to 40 per cent and the best power stations about 60 per cent; it is also why heat flows from
hot to cold and never back on its own, and why a perpetual-motion machine is impossible.
\item \textbf{Entropy gives time a direction.} The same law says that the entropy of an isolated system
never decreases, which is why heat spreads out, why gases mix but never spontaneously un-mix, and why
the future looks different from the past.
\end{itemize}
\end{motivbox}

% ================================================================
\section{Systems, Equilibrium and the Zeroth Law}\label{sec:thermo-zeroth}

Before stating any law we need words for what the law is about. Thermodynamics always divides the world
in two: the part we study and everything else. The laws then speak about what crosses the boundary
between them.

\begin{defbox}{System, surroundings, equilibrium}
The \textbf{system} is the body or region under study (the gas in a cylinder, the water in a kettle);
everything that can exchange energy with it is the \textbf{surroundings}. The boundary may let energy
through as heat (a \emph{diathermic} wall, such as thin metal) or block heat (an \emph{adiabatic} wall,
an ideal insulator), and it may be fixed or movable (a piston). A system is in \textbf{thermodynamic
equilibrium} when its macroscopic properties, such as pressure, volume and temperature, are uniform
throughout and do not change with time. Two systems in contact through a diathermic wall are in
\textbf{thermal equilibrium} with each other when no heat flows between them.
\end{defbox}

\subsection{Thermal equilibrium and the zeroth law}

\chref{ch:thermal} introduced temperature as the property that decides which way heat flows, used a thermometer to measure it, and named in passing the assumption hidden in that procedure. It deserves a formal statement here, at the head of the laws, because nothing else in thermodynamics would make sense without it.

\begin{defbox}{The zeroth law of thermodynamics}
If body $A$ is in thermal equilibrium with body $C$, and body $B$ is separately in thermal equilibrium
with $C$, then $A$ and $B$ are in thermal equilibrium with each other.
\end{defbox}

\why{This sounds too obvious to call a law until you notice that it is what makes a \emph{thermometer}
possible. The thermometer is body $C$: it comes to equilibrium with the patient, and its reading can then
be compared with the reading it gives for a bath or a bottle the patient never touches. The law
guarantees that ``same temperature'' is a consistent label: all bodies in equilibrium with one another
share a single number, the temperature. It is a fact of experience, not a theorem; nothing in mechanics
implies it. It was called ``zeroth'' because its logical necessity was recognised only after the first
and second laws already had their numbers.}

\subsection{Internal energy, heat and work}

Three energies appear in every thermodynamic argument, and most errors in the subject come from
confusing them. One is a store; the other two are ways of filling or emptying it.

\begin{defbox}{Internal energy, heat and work}
\begin{itemize}
\item \textbf{Internal energy} $U$ (\si{J}, a scalar) is the total energy of the molecules of the system
      measured in the frame of its centre of mass (CM): their random kinetic energy plus the potential
      energy of their mutual forces. It excludes the kinetic energy of the system moving as a whole and
      its potential energy in external fields. $U$ is a \textbf{state variable}: its value is fixed by the
      present equilibrium state, for example by $(P,V,T)$, and not by the history that led there.
\item \textbf{Heat} $Q$ (\si{J}) is energy transferred between the system and its surroundings
      \emph{because of a temperature difference} between them.
\item \textbf{Work} $W$ (\si{J}) is energy transferred by a macroscopic force acting through a macroscopic
      displacement: for a gas, by a moving piston.
\end{itemize}
$Q$ and $W$ are \emph{not} state variables. A system does not ``contain'' heat or work; they describe
energy in transit, and how much is transferred depends on the path taken between two states.
\end{defbox}

\begin{remarkbox}{Heat is not a substance}
Until the 1840s most scientists pictured heat as an invisible fluid, ``caloric'', that flowed from hot
bodies to cold ones and was conserved. Rumford's observation that boring a cannon produced heat without limit, and Joule's careful measurements that a falling weight stirring water warms it by an amount proportional to the work done (\chref{ch:thermal}, Section~\ref{sec:calorimetry}), ended that picture: \SI{4.186}{J} of work produces the same warming as one calorie of heat. Heat and work are two routes for the same thing, energy, which is exactly what
the first law will say. The phrase ``heat content of a body'' is a relic of the caloric theory; this book
avoids it.
\end{remarkbox}

\subsection{Work done by an expanding gas}

Of the two transfer modes, work is the one we can compute from mechanics alone. Doing so now gives the
tool used in every process of this chapter, and it shows directly why work depends on the path.

\begin{derivbox}{Work done by a gas on a moving piston}\label{der:thermo-pdv}
\stepn{1} A gas at pressure $P$ pushes on a piston of face area $A$ with a force $F=PA$
(Figure~\ref{fig:thermo-piston}).

\stepn{2} If the piston moves out a small distance $ds$, slowly enough that the pressure stays uniform
and equal to $P$, the work done \emph{by} the gas is
\[ dW = F\,ds = PA\,ds . \]

\stepn{3} But $A\,ds$ is exactly the volume swept out, $dV$. So $dW=P\,dV$, and over a finite change from
volume $V_1$ to $V_2$,
\[
\boxed{\;W=\int_{V_1}^{V_2} P\,dV\;}
\]

\stepn{4} \textbf{Geometric reading.} On a graph of $P$ against $V$, the integral is the \textbf{area under
the curve} between $V_1$ and $V_2$ (the same reading of an integral as the area under a force--distance
graph in \chref{ch:energy}). Expansion ($dV>0$) gives positive work; compression gives negative work.
\end{derivbox}

\begin{figure}[htbp]
\centering
\tikzfig{thermo_piston}
\figcaption{fig:thermo-piston}{The gas pushes the piston outwards with force $PA$. As the piston moves
from its present position to the faded one, a distance $ds$, it sweeps out the orange slice of volume
$dV=A\,ds$, and the gas does work $PA\,ds=P\,dV$. Nothing here depends on the shape of the cylinder, so
$dW=P\,dV$ holds for any container whose walls move.}
\end{figure}

\begin{trapbox}{whose pressure goes into $\int P\,dV$}
The derivation used the pressure of the gas, and that is legitimate only if the gas has a single, uniform
pressure throughout, that is, if the piston moves slowly (we shall call such a change quasi-static in
Section~\ref{sec:quasistatic}). In a sudden expansion the gas near the piston is thinner than the gas far
from it, and there is no single $P$ to use. What is always true is that the work done by the gas equals
the work done \emph{against} the external pressure on the piston, $W=\int P_{\text{ext}}\,dV$. A gas
expanding into a vacuum pushes on nothing, so it does no work at all, however much its volume grows.
\end{trapbox}

\subsubsection{Why work depends on the path: a concrete demonstration}

The claim that $W$ depends on the path deserves more than assertion. Take a gas between the same two
states by two different routes and compare the numbers.

\begin{figure}[htbp]
\centering
\tikzfig{thermo_twopaths}
\figcaption{fig:thermo-twopaths}{Two routes between the same states $i$ and $f$. The work done by the gas
is the area under each path. Path~1 expands while the pressure is high and encloses both shaded
regions; path~2 drops the pressure first and encloses only the lower one. The change in internal energy,
a state variable, is the same for both.}
\end{figure}

\begin{derivbox}{Same endpoints, different work}\label{der:thermo-twopaths}
Take $i=(P,V)=(\SI{2e5}{Pa},\,\SI{1e-3}{m^3})$ and $f=(\SI{1e5}{Pa},\,\SI{2e-3}{m^3})$, the two corners
in Figure~\ref{fig:thermo-twopaths}.

\stepn{1} \textbf{Path 1} (expand first, at high pressure): the constant-pressure leg does
$W=P\,\Delta V=(\num{2e5})(\num{1e-3})=\SI{200}{J}$; the constant-volume leg does none, since the piston
does not move. Total $W_1=\SI{200}{J}$.

\stepn{2} \textbf{Path 2} (drop the pressure first, then expand): now the expansion happens at the lower
pressure, $W=(\num{1e5})(\num{1e-3})=\SI{100}{J}$. Total $W_2=\SI{100}{J}$.

\stepn{3} Same start, same finish, yet $W_1\neq W_2$. The internal energy change $\Delta U$ is fixed by the
endpoints alone, so (as the first law in the next section will show) the \emph{heats} supplied along the
two paths must also differ, by exactly \SI{100}{J}.
\end{derivbox}

\subsection{The sign convention}

Heat and work can each flow either way, so we must fix signs once and keep them. This book follows the
convention of the NCERT textbooks and of most physics courses.

\begin{keybox}{Sign convention used throughout this book}\label{key:thermo-signs}
\begin{itemize}
\item $Q>0$ when heat flows \emph{into} the system; $Q<0$ when the system gives heat out.
\item $W>0$ when work is done \emph{by} the system on its surroundings (expansion); $W<0$ when work is done
      \emph{on} the system (compression).
\end{itemize}
With these signs the first law reads $Q=\Delta U+W$. Many chemistry books, and some physics books,
instead let $W$ stand for the work done \emph{on} the gas; their first law is $\Delta U=Q+W$. The physics is
identical, only the letter's meaning differs. Check the convention before borrowing any formula.
\end{keybox}

% ================================================================
\section{The First Law and the Specific Heats of a Gas}\label{sec:firstlaw}

We now have the three energies, internal energy $U$, heat $Q$ and work $W$, and a sign for each. The
first law states how they balance. Energy is never created or destroyed, only moved between the store
($U$) and the two doors through which it enters or leaves (heat and work). Everything in the rest of the
chapter is this one equation applied to different situations.

\subsection{The first law}

\begin{keybox}{First law of thermodynamics}\label{key:firstlaw}
For any process taking a system from one equilibrium state to another,
\[
\boxed{\;Q=\Delta U+W\;}
\qquad\text{that is,}\qquad
\text{heat in} = \text{rise in internal energy}+\text{work done by the system}.
\]
For a small step of a quasi-static process, $\delta Q=dU+P\,dV$. Here $dU$ is a small change in the
state variable $U$, while $\delta Q$ and $P\,dV$ are small \emph{amounts} transferred; the different
symbol is a reminder that there is no ``$Q$ of the state'' whose change $\delta Q$ is.
\end{keybox}

We write $Q$ and $W$, not $\Delta Q$ and $\Delta W$, for the same reason: a $\Delta$ means ``final value
minus initial value'' of something the system has, and the system has no value of heat or work.

\inwords{Feed a gas some heat. It can bank the energy (raise $U$, which for a gas means getting hotter) or
spend it (push the piston, doing work), and the first law says the books balance exactly. Every process
in Section~\ref{sec:processes} is a different split between banking and spending.}

\why{The first law is why a car engine, a power station and a pressure cooker can be designed on paper.
Burning petrol releases heat; the engine's job is to turn as much of it as possible into work rather than
into internal energy of the engine block. And the law's hard limit, that you cannot get out more work than
the energy you put in, is the reason a ``free-energy'' machine (a \emph{perpetual-motion machine of the
first kind}) is impossible, however clever the mechanism.}

The law is not a definition in disguise. Its content is that $U$, defined as heat in minus work out,
\emph{is} a state variable: measure $Q-W$ along any path between two given states and you get the same
number. Figure~\ref{fig:thermo-twopaths} shows the content at work: $W$ differed by \SI{100}{J} between the
two paths, so the heats must differ by the same \SI{100}{J}, and they do.

\begin{exbox}{Boiling water: where does the latent heat go?}\label{ex:thermo-boiling}
At atmospheric pressure ($\SI{1.013e5}{Pa}$), \SI{1.00}{g} of water at \SI{100}{\celsius} occupies
\SI{1.00}{cm^3}; as steam at the same temperature it occupies \SI{1671}{cm^3}. The latent heat of
vaporisation (\chref{ch:thermal}, Section~\ref{sec:latent}) is $L_v=\SI{2256}{J/g}$. How much of this heat does work on the
atmosphere, and how much raises the internal energy?

\stepn{1} \textbf{Heat in.} $Q=mL_v=(1.00)(2256)=\SI{2256}{J}$, positive because heat enters the water.

\stepn{2} \textbf{Work done by the water.} The steam pushes back the atmosphere at constant pressure:
$W=P\,\Delta V=(\num{1.013e5})(1671-1)\times10^{-6}=\SI{169}{J}$.

\stepn{3} \textbf{Internal energy.} $\Delta U=Q-W=2256-169=\SI{2087}{J}$.

\stepn{4} \textbf{Read the result.} About 93 per cent of the latent heat goes into pulling the molecules
apart against their mutual attraction (a rise in intermolecular potential energy at constant temperature);
only 7 per cent pays for pushing the atmosphere out of the way. Water vapour is not an ideal gas here: its
$U$ changed at constant $T$, which the next caution says an ideal gas cannot do.
\ans{$W\approx\SI{169}{J}$, $\Delta U\approx\SI{2087}{J}$}
\end{exbox}

\begin{trapbox}{for an ideal gas, $U$ depends on $T$ alone}
In an ideal gas the molecules exert no forces on one another except in collisions, so there is no
intermolecular potential energy and $U$ is purely the kinetic energy of molecular motion.
\chref{ch:kinetic} shows that this kinetic energy is fixed by the temperature alone; experimentally,
Joule found that a real gas expanding freely into a vacuum changes its temperature hardly at all, and the
ideal gas is the limit in which the change is exactly zero. Two consequences are used constantly:
\begin{itemize}
\item an isothermal process ($T$ constant) has $\Delta U=0$, whatever happens to $P$ and $V$;
\item a free expansion into a vacuum ($Q=0$, $W=0$) has $\Delta U=0$, so $T$ is unchanged.
\end{itemize}
For real gases, liquids and solids, $U$ depends on volume too, as the boiling water above shows.
\end{trapbox}

\subsection{Why a gas has two specific heats, and why \texorpdfstring{$C_p>C_v$}{Cp > Cv}}

\chref{ch:thermal} (Section~\ref{sec:calorimetry}) defined the specific heat $s$ (per kilogram) and the molar heat capacity $C$ (per
mole) as the heat needed per unit temperature rise. For a solid or liquid that is one number, because
its volume barely changes when heated. For a gas the answer depends on what is held fixed while the heat
goes in.

Heat a gas in a sealed rigid box, and every joule goes into $U$. Heat it under a freely sliding piston, at
constant pressure, and it expands as it warms, so some of the joules leave again as work and \emph{more}
heat is needed for the same temperature rise. Hence $C_p>C_v$, before any algebra.

\why{This is not hair-splitting. A parcel of air rising in the atmosphere expands against the air around
it, a rigid gas cylinder left in the sun does not, and the two warm or cool at different rates for the
same heat. The ratio of the two heat capacities also sets how fast sound travels through a gas, and
whether a diesel engine's compression stroke gets hot enough to ignite the fuel.}

\begin{defbox}{The two molar heat capacities of a gas}
For $n$ moles of gas taking in a small amount of heat $\delta Q$ while its temperature rises by $dT$,
\[
C_v=\frac{1}{n}\left(\frac{\delta Q}{dT}\right)_{V\ \text{constant}},
\qquad
C_p=\frac{1}{n}\left(\frac{\delta Q}{dT}\right)_{P\ \text{constant}},
\]
both in \si{J.mol^{-1}.K^{-1}}. So $Q=nC_v\,\Delta T$ when the volume is clamped, and $Q=nC_p\,\Delta T$
when the pressure is held fixed. The subscript names what was held fixed while the heat went in; for a gas,
a heat capacity without that information is meaningless.
\end{defbox}

\begin{derivbox}{Mayer's relation, \texorpdfstring{$C_p-C_v=R$}{Cp - Cv = R}}\label{der:mayer}
Take $n$ moles of ideal gas through the same small temperature rise $dT$, twice.

\stepn{1} \textbf{At constant volume.} The piston does not move, so $W=0$ and the first law collapses to
$\delta Q=dU$. By the definition of $C_v$,
\[ dU=nC_v\,dT . \]
Because $U$ of an ideal gas depends only on $T$, \emph{this expression for $dU$ holds in every process},
not only at constant volume: any two states at the same pair of temperatures have the same $\Delta U$. We
use this at once.

\stepn{2} \textbf{At constant pressure.} Now the gas expands, and the first law reads
\[ nC_p\,dT=dU+P\,dV=nC_v\,dT+P\,dV . \]

\stepn{3} Evaluate $P\,dV$ from the ideal-gas equation. At constant pressure, a small change of $PV=nRT$
gives
\[ P\,dV=nR\,dT . \]

\stepn{4} Substitute and divide by $n\,dT$:
\[ C_p=C_v+R \qquad\Longrightarrow\qquad \boxed{\;C_p-C_v=R\;} \]

\stepn{5} \textbf{Read the result.} The difference is exactly the work, $R$ per mole per kelvin, that the gas
does while expanding at constant pressure. $R$ here is not an abstract constant; it is the rate at which
heat leaks out as work. For liquids and solids, which barely expand, $C_p$ and $C_v$ are nearly equal,
and for real gases at high density the relation is only approximate.
\end{derivbox}

\begin{keybox}{The ratio \texorpdfstring{$\gamma$}{gamma}, and \texorpdfstring{$C_v$}{Cv} in terms of it}\label{key:gamma}
The ratio of the two heat capacities is
\[ \gamma\equiv\frac{C_p}{C_v}>1\qquad\text{always, since } C_p=C_v+R. \]
Dividing Mayer's relation by $C_v$ gives $\gamma-1=R/C_v$, so
\[
\boxed{\;C_v=\frac{R}{\gamma-1}\;},\qquad C_p=\frac{\gamma R}{\gamma-1}.
\]
These let every later formula trade $C_v$ for $\gamma$. \chref{ch:kinetic} predicts, from molecular
structure alone, $C_v=\tfrac32R$ and $\gamma=\tfrac53$ for a monatomic gas, and $C_v=\tfrac52R$,
$\gamma=\tfrac75$ for a diatomic gas near room temperature; air, which is mostly diatomic, has
$\gamma=1.40$. Until then we take these values as given.
\end{keybox}

% ================================================================
\section{Thermodynamic Processes}\label{sec:processes}

The first law says energy is conserved, but not yet how a gas behaves as it is compressed in an engine
cylinder, expanded through a nozzle or heated in a sealed can. That depends on \emph{how} the change is
made: which variable is held fixed, and whether heat can get in. There are four standard ways, at constant
temperature (isothermal), with no heat exchanged (adiabatic), at constant volume (isochoric) and at
constant pressure (isobaric), and they are the building blocks of every engine cycle. This section first
says what may be drawn on a $P$--$V$ diagram, then treats all four processes as members of one family.

\subsection{State variables and quasi-static change}\label{sec:quasistatic}

\begin{defbox}{State variables and the equation of state}
\begin{itemize}
\item A \textbf{state variable} is any quantity fixed by the system's present equilibrium state, with no
      memory of its history: $P$, $V$, $T$ and $U$ all qualify; $Q$ and $W$ do not.
\item State variables come in two kinds. \textbf{Extensive} ones scale with the amount of substance ($V$,
      $U$, $n$); \textbf{intensive} ones do not ($P$, $T$). Divide a gas in two with a partition, and each
      half has half the volume but the same pressure and temperature.
\item The variables are linked by an \textbf{equation of state}. For $n$ moles of ideal gas it is
      $PV=nRT$, so fixing any two of $(P,V,T)$ fixes the third. That is why a single point on the
      $P$--$V$ plane specifies the state completely.
\end{itemize}
\end{defbox}

\begin{defbox}{Quasi-static process}
An equation of state holds only in equilibrium. A gas in the middle of an explosion has no single
pressure or temperature, and so no point on the diagram. A \textbf{quasi-static} process is carried out
so slowly that the system passes through a continuous chain of equilibrium states, and \emph{that} is
what a curve on the $P$--$V$ diagram depicts. Every process drawn as a curve in this chapter is
quasi-static by assumption; Section~\ref{sec:reversible} adds the stricter idea of a \emph{reversible}
process.
\end{defbox}

``Slowly'' means slowly compared with the time the gas takes to even out its pressure, which for a
litre of gas is about the time sound takes to cross it, well under a millisecond. So the compression
stroke of a car engine, a few milliseconds long, is close to quasi-static as far as pressure goes, even
though it is far too fast for much heat to leak out; we shall use both facts.

A \emph{process} is then a path on the $P$--$V$ diagram. For each one we ask the same three questions:
what is $W$, what is $\Delta U$, and what does the first law then say about $Q$?

\subsection{One relation for all four: the polytropic process}\label{sec:polytropic}

The four textbook processes look like four separate topics. They are not: each is a special case of a
single relation, and deriving the general case once makes the special cases almost free.

\begin{defbox}{Polytropic process}
A \textbf{polytropic} process is a quasi-static process of an ideal gas that obeys
\[
\boxed{\;PV^{x}=\text{constant}\;}
\]
for some fixed number $x$, the \textbf{polytropic index} (a pure number, which may be any real value).
The four standard processes are four choices of $x$:
\begin{center}
\begin{tabular}{@{}clll@{}}
\toprule
$x$ & Relation & Process & Held fixed\\
\midrule
$0$ & $P=$ constant & isobaric & pressure $P$\\
$1$ & $PV=$ constant & isothermal & temperature $T$\\
$\gamma$ & $PV^\gamma=$ constant & adiabatic & no heat, $Q=0$\\
$\infty$ & $V=$ constant & isochoric & volume $V$\\
\bottomrule
\end{tabular}
\end{center}
\end{defbox}

The rows for $x=0$ and $x=1$ are immediate ($PV=nRT$ is constant exactly when $T$ is). The row for the
adiabatic process is a result, derived in Section~\ref{sec:adiabatic}. The last row needs a word: rewrite
$PV^x=$ constant as $P^{1/x}V=$ constant; as $x\to\infty$, $1/x\to0$, $P^{1/x}\to1$ for any fixed $P$, and
the relation becomes $V=$ constant. The ``infinity'' is an honest limit, not a trick.

\inwords{Read the index as a dial. Turn it from $0$ up through $1$ and $\gamma$ towards $\infty$, and the
$P$--$V$ curve through a given state tilts continuously from flat (isobaric) to vertical (isochoric),
passing through the isotherm and the adiabat on the way. Everything in this section is one journey along
that dial.}

\begin{trapbox}{two different $n$'s in other books}
Many books write the polytropic law as $PV^n=$ constant and then use $n$ for the number of moles in
$PV=nRT$ on the same page, so a formula such as $W=nR(T_1-T_2)/(n-1)$ has two different $n$'s in it,
the top one moles and the bottom one the index. This book keeps $n$ for moles and writes the index as
$x$, so the same formula reads $W=nR(T_1-T_2)/(x-1)$.
\end{trapbox}

\subsubsection{The shapes of the curves}

\begin{figure}[htbp]
\centering
\tikzfig{thermo_polytropes}
\figcaption{fig:thermo-polytropes}{The polytropic family through one state $(V_0,P_0)$, drawn for a
monatomic gas ($\gamma=\tfrac53$). The slope at the common point is $-xP_0/V_0$, so the curves steepen in
order of $x$: flat for the isobar, vertical for the isochore. The adiabat is the steeper of the two
falling curves, by the factor $\gamma$, and it lies below the isotherm on the expansion side.}
\end{figure}

\begin{derivbox}{The slope of a polytropic curve}\label{der:polyslope}
\stepn{1} Differentiate $PV^x=$ constant with respect to $V$, using the product rule:
\[ \frac{dP}{dV}\,V^x+P\,xV^{x-1}=0 . \]

\stepn{2} Solve for the slope:
\[ \frac{dP}{dV}=-\,\frac{xPV^{x-1}}{V^x}=-\,x\,\frac{P}{V} . \]

\stepn{3} At a shared point $(P,V)$ the slope is $-x$ times $P/V$: the larger $x$, the steeper the fall
(Figure~\ref{fig:thermo-polytropes}). In particular the adiabat's slope is exactly $\gamma$ times the
isotherm's, the most useful single fact for sketching these diagrams. For negative $x$ the slope is
positive: pressure and volume rise together, as in a process with $P\propto V$ ($x=-1$).
\end{derivbox}

\subsection{The general work integral}

We now compute the work once, for any index, and collect the special cases afterwards.

\begin{derivbox}{Work in a polytropic process (\texorpdfstring{$x\neq1$}{x not 1})}\label{der:polywork}
\stepn{1} Along the path $PV^x=P_1V_1^{\,x}$, so $P=P_1V_1^{\,x}V^{-x}$. Insert into
$W=\int_{V_1}^{V_2}P\,dV$:
\[ W=P_1V_1^{\,x}\int_{V_1}^{V_2}V^{-x}\,dV . \]

\stepn{2} Integrate the power, which is allowed whenever $x\neq1$, so that the exponent $1-x$ is not zero:
\[ W=P_1V_1^{\,x}\,\frac{V_2^{\,1-x}-V_1^{\,1-x}}{1-x} . \]

\stepn{3} \textbf{Tidy up.} The constant may be written either as $P_1V_1^{\,x}$ or as $P_2V_2^{\,x}$, so
$P_1V_1^{\,x}V_2^{\,1-x}=P_2V_2^{\,x}V_2^{\,1-x}=P_2V_2$ and $P_1V_1^{\,x}V_1^{\,1-x}=P_1V_1$. Hence
\[ W=\frac{P_2V_2-P_1V_1}{1-x}=\boxed{\;\frac{P_1V_1-P_2V_2}{\,x-1\,}\;} \]

\stepn{4} \textbf{Temperature form.} Since $PV=nRT$ at every point, $P_1V_1=nRT_1$ and $P_2V_2=nRT_2$:
\[ \boxed{\;W=\frac{nR\,(T_1-T_2)}{\,x-1\,}\;} \]

\stepn{5} \textbf{Check.} For $x=0$ (constant pressure) this gives $W=nR(T_2-T_1)=P(V_2-V_1)$, the
rectangle under a horizontal line, as it must. As $x\to1$ the formula becomes $0/0$; that case needs its
own integral, done in Section~\ref{sec:isothermal}.
\end{derivbox}

\subsection{The general heat capacity}

Here is the payoff that the four-separate-processes approach hides completely: \emph{every} path has its
own molar heat capacity, and one formula gives all of them.

\begin{derivbox}{The molar heat capacity along a polytropic path}\label{der:polyC}
\stepn{1} Define $C_x$ as the molar heat capacity \emph{along the path}: the heat taken in is
$Q=nC_x\,\Delta T$ for a temperature change $\Delta T=T_2-T_1$.

\stepn{2} Use the first law with two results already in hand: $\Delta U=nC_v\,\Delta T$ (true for every
process of an ideal gas, Derivation~\ref{der:mayer}), and the work from Derivation~\ref{der:polywork},
\[
W=\frac{nR(T_1-T_2)}{x-1}=\frac{nR\,\Delta T}{1-x}.
\]

\stepn{3} Assemble $Q=\Delta U+W$ and divide by $n\,\Delta T$:
\[
C_x\,\Delta T=C_v\,\Delta T+\frac{R\,\Delta T}{1-x}
\quad\Longrightarrow\quad
\boxed{\;C_x=C_v+\frac{R}{1-x}\;}
\]

\stepn{4} With $C_v=R/(\gamma-1)$ (Result~\ref{key:gamma}) this can be written in terms of $\gamma$ and $x$
alone:
\[
C_x=R\left(\frac{1}{\gamma-1}+\frac{1}{1-x}\right)=R\,\frac{\gamma-x}{(\gamma-1)(1-x)} .
\]
\end{derivbox}

\begin{keybox}{Every process's heat capacity from one formula}\label{key:polyC}
\[
C_x=C_v+\frac{R}{1-x}=R\,\frac{\gamma-x}{(\gamma-1)(1-x)}
\]
\begin{center}
\begin{tabular}{@{}cllc@{}}
\toprule
$x$ & Process & $C_x$ & Diatomic gas\\
\midrule
$0$ & isobaric & $C_v+R=C_p$ & $\tfrac72R$\\
$1$ & isothermal & infinite & infinite\\
$\gamma$ & adiabatic & $0$ & $0$\\
$\infty$ & isochoric & $C_v$ & $\tfrac52R$\\
\bottomrule
\end{tabular}
\end{center}
\end{keybox}

\inwords{The heat capacity answers ``how much heat per kelvin?'', and each row tells a story. Isochoric: no
work, so all the heat becomes temperature, and the capacity is $C_v$. Isobaric: the gas also does work as
it expands, so more heat is needed per kelvin, $C_p>C_v$. Isothermal: heat flows in but the temperature
does not move, so the heat per kelvin is infinite. Adiabatic: no heat at all, yet the temperature
changes, so the capacity is zero.}

\begin{trapbox}{a heat capacity can be zero, infinite or negative}
For a gas, ``the specific heat'' is not one number; it depends on the path. Stranger still, for
$1<x<\gamma$ the formula gives a \textbf{negative} $C_x$: the numerator $\gamma-x$ is positive and the
factor $1-x$ is negative. Such a gas \emph{gives out} heat while its temperature \emph{rises} (in a
compression), or takes heat in while it cools (in an expansion). There is no paradox: in a compression the
work done on the gas exceeds the rise in $U$, and the surplus leaves as heat.
Example~\ref{ex:thermo-polytropic} shows the numbers.
\end{trapbox}

\subsection{The four special cases}

Three of the four cases are now substitutions; the fourth, the isotherm, is the one value of $x$ that the
general integral skipped.

\subsubsection{Isothermal process (\texorpdfstring{$x=1$}{x = 1})}\label{sec:isothermal}

At $x=1$ the exponent $1-x$ in the work integral vanishes and the power rule fails; the integral becomes
a logarithm instead.

\begin{derivbox}{Isothermal work}\label{der:isowork}
\stepn{1} With $T$ fixed, $P=nRT/V$, so
\[ W=\int_{V_1}^{V_2}\frac{nRT}{V}\,dV=nRT\int_{V_1}^{V_2}\frac{dV}{V} . \]

\stepn{2} The integral of $1/V$ is $\ln V$ (\chref{ch:toolkit}, Section~\ref{sec:explog}), so
\[ \boxed{\;W=nRT\,\ln\frac{V_2}{V_1}=nRT\,\ln\frac{P_1}{P_2}\;} \]
using $P_1V_1=P_2V_2$ for the second form. The logarithm carries the sign: positive for expansion,
negative for compression.

\stepn{3} Since $\Delta T=0$, $\Delta U=0$, and the first law gives $Q=W$: every joule of heat that enters
leaves as work. Consistently, $C_1=Q/(n\,\Delta T)$ is infinite, because $Q\neq0$ while $\Delta T=0$.
\end{derivbox}

An isothermal process needs the gas to stay in good thermal contact with a large body at fixed
temperature (a \emph{reservoir}, defined in Section~\ref{sec:heatengine}) and to change slowly enough for
heat to flow across with only an infinitesimal temperature difference. It is an idealisation, approached
by slow changes in thin-walled containers.

\subsubsection{Adiabatic process (\texorpdfstring{$x=\gamma$}{x = gamma})}\label{sec:adiabatic}

The other cases fix $x$ by holding $P$, $T$ or $V$ constant. The adiabatic case is different: nothing is
held constant except that no heat crosses the boundary, and the index $\gamma$ is \emph{forced} by that
demand.

\begin{derivbox}{\texorpdfstring{$Q=0$}{Q = 0} forces \texorpdfstring{$PV^\gamma=$}{PV gamma =} constant}\label{der:adiabat}
\stepn{1} Set $\delta Q=0$ in the first law for a small quasi-static step: $0=nC_v\,dT+P\,dV$.

\stepn{2} A small change of $PV=nRT$ gives $P\,dV+V\,dP=nR\,dT$, so $n\,dT=(P\,dV+V\,dP)/R$.

\stepn{3} Substitute and multiply by $R$: $\;C_v(P\,dV+V\,dP)+RP\,dV=0$.

\stepn{4} Replace $R$ by $C_p-C_v$ (Mayer) and collect: $\;C_vV\,dP+C_pP\,dV=0$.

\stepn{5} Divide by $C_vPV$ and write $\gamma=C_p/C_v$:
\[ \frac{dP}{P}+\gamma\,\frac{dV}{V}=0 . \]

\stepn{6} Integrate (each term gives a logarithm, \chref{ch:toolkit}, Section~\ref{sec:explog}):
$\ln P+\gamma\ln V=$ constant, that is,
\[ \boxed{\;PV^\gamma=\text{constant}\;} \]
the polytropic law with $x=\gamma$. Using $P=nRT/V$ gives the equivalent forms $TV^{\gamma-1}=$ constant
and $P^{1-\gamma}T^\gamma=$ constant.
\end{derivbox}

\begin{keybox}{Adiabatic work, two ways}\label{key:adiawork}
Setting $x=\gamma$ in the polytropic work gives
\[
W=\frac{nR(T_1-T_2)}{\gamma-1}=\frac{P_1V_1-P_2V_2}{\gamma-1},
\]
and the same result follows in one line from the first law: with $Q=0$,
$W=-\Delta U=nC_v(T_1-T_2)=nR(T_1-T_2)/(\gamma-1)$. Since an expansion does positive work, it must have
$T_2<T_1$: \textbf{an adiabatic expansion cools the gas, and an adiabatic compression heats it}. Air rising
in the atmosphere expands and cools in this way until its water vapour condenses into cloud; a bicycle
pump grows warm for the reverse reason.
\end{keybox}

\begin{trapbox}{``adiabatic'' does not mean ``constant temperature''}
It means \emph{no heat flow}. The temperature changes precisely \emph{because} no heat can flow: the work
must be paid for out of $U$. Isothermal and adiabatic are opposite extremes: very slow with perfect thermal
contact ($x=1$), against perfectly insulated or too fast for heat to move ($x=\gamma$). A diesel engine
compresses air by a factor of about 20 in a few milliseconds, close to adiabatic; from \SI{300}{K},
$T_2=300\times20^{0.4}\approx\SI{990}{K}$, hot enough to ignite the fuel without a spark. Real
processes lie between the extremes, with $1<x<\gamma$.
\end{trapbox}

\subsubsection{Isobaric (\texorpdfstring{$x=0$}{x = 0}) and isochoric (\texorpdfstring{$x\to\infty$}{x to infinity}) processes}

\begin{derivbox}{Isobaric process (\texorpdfstring{$P$}{P} constant)}\label{der:isobaric}
\stepn{1} \textbf{Work.} $P$ comes straight out of the integral:
\[ W=\int_{V_1}^{V_2}P\,dV=P(V_2-V_1)=nR\,\Delta T, \]
the last step from $PV=nRT$ at constant $P$.

\stepn{2} \textbf{Heat.} At $x=0$, $C_0=C_v+R=C_p$, as it must be, so $Q=nC_p\,\Delta T$.

\stepn{3} \textbf{First-law check.} $\Delta U=nC_v\,\Delta T$ and $W=nR\,\Delta T$, so
$Q=\Delta U+W=n(C_v+R)\,\Delta T=nC_p\,\Delta T$, in agreement with Step~2.
\end{derivbox}

\begin{derivbox}{Isochoric process (\texorpdfstring{$V$}{V} constant)}\label{der:isochoric}
\stepn{1} \textbf{Work.} The volume never changes, so $dV=0$ and $W=\int P\,dV=0$: the walls do not move.

\stepn{2} \textbf{Heat capacity.} As $x\to\infty$, $R/(1-x)\to0$ and $C_x\to C_v$, the constant-volume
capacity, as its name requires.

\stepn{3} \textbf{First law.} With $W=0$, $Q=\Delta U=nC_v\,\Delta T$: every joule of heat goes into
internal energy, and so into temperature.
\end{derivbox}

\subsection{The four processes in one table}

\begin{tablebox}{All four processes as one family (\texorpdfstring{$n$}{n} moles of ideal gas)}
\begin{center}
\renewcommand{\arraystretch}{1.35}
\begin{tabular}{@{}cllllc@{}}
\toprule
$x$ & Process & Path & Work $W$ & $\Delta U$ & $C_x$\\
\midrule
$0$ & isobaric & $P=$ constant & $P\,\Delta V=nR\,\Delta T$ & $nC_v\,\Delta T$ & $C_p$\\
$1$ & isothermal & $PV=$ constant & $nRT\ln\dfrac{V_2}{V_1}$ & $0$ & infinite\\
$\gamma$ & adiabatic & $PV^\gamma=$ constant & $\dfrac{nR(T_1-T_2)}{\gamma-1}$ & $-W$ & $0$\\
$\infty$ & isochoric & $V=$ constant & $0$ & $nC_v\,\Delta T$ & $C_v$\\
\midrule
$x$ & polytropic & $PV^x=$ constant & $\dfrac{nR(T_1-T_2)}{x-1}$ & $nC_v\,\Delta T$ & $C_v+\dfrac{R}{1-x}$\\
\bottomrule
\end{tabular}\label{tab:processes}
\end{center}
Read the bottom row and set $x$ to $0$, $1$, $\gamma$ and $\infty$ to recover the rows above it. One
entry never changes: $\Delta U=nC_v\,\Delta T$ in \emph{every} row (in the adiabatic row,
$-W=nC_v\,\Delta T$ too), because $U$ of an ideal gas depends only on $T$. All the variety lies in how the
heat divides between $\Delta U$ and $W$.
\end{tablebox}

\begin{exbox}{One gas, two routes, different work}\label{ex:thermo-isoadia}
Two moles of a monatomic ideal gas ($\gamma=\tfrac53$) at \SI{300}{K} double their volume (a)
isothermally, (b) adiabatically. Find the work done by the gas in each case, and the final temperature in
(b). Take $R=\SI{8.314}{J.mol^{-1}.K^{-1}}$.

\stepn{1} \textbf{(a) Isothermal.} $W=nRT\ln(V_2/V_1)=2(8.314)(300)\ln2\approx\SI{3458}{J}$. The temperature
stays at \SI{300}{K}, so $\Delta U=0$ and $Q=W=\SI{3458}{J}$.

\stepn{2} \textbf{(b) Adiabatic.} Use $TV^{\gamma-1}=$ constant with $\gamma-1=\tfrac23$:
\[ T_2=300\times2^{-2/3}\approx\SI{189.0}{K}. \]

\stepn{3} Work from the temperature drop:
\[ W=\frac{nR(T_1-T_2)}{\gamma-1}=\frac{2(8.314)(300-189.0)}{2/3}\approx\SI{2769}{J}. \]

\stepn{4} \textbf{Compare.} The adiabatic route does less work ($2769<3458$), as the steeper adiabat of
Figure~\ref{fig:thermo-isoadia} predicts: it encloses less area. All of it is paid for out of internal
energy, which is why the gas cools by \SI{111}{K}; on the isothermal route the surroundings keep topping
up the gas with heat.
\ans{(a) \SI{3458}{J} at \SI{300}{K}; (b) \SI{2769}{J}, final $T\approx\SI{189}{K}$}
\end{exbox}

\begin{figure}[htbp]
\centering
\tikzfig{thermo_isoadia}
\figcaption{fig:thermo-isoadia}{Example~\ref{ex:thermo-isoadia} drawn to scale. Both curves start at the
same state and end at twice the volume. The teal area under the adiabat is its work, \SI{2769}{J}; the
isotherm adds the pale blue sliver, \SI{689}{J}, for a total of \SI{3458}{J}. The adiabat ends at a lower
pressure because the gas has cooled.}
\end{figure}

\begin{exbox}{A polytropic compression, and a negative heat capacity}\label{ex:thermo-polytropic}
One mole of a diatomic ideal gas ($\gamma=1.4$, $C_v=\tfrac52R$) is compressed along $PV^{1.3}=$ constant
from $V_1=\SI{2}{L}$ to $V_2=\SI{1}{L}$, starting at $T_1=\SI{300}{K}$. Find $T_2$, $W$, $\Delta U$, $Q$ and
the heat capacity $C_x$, and comment on the signs.

\stepn{1} \textbf{Final temperature.} With $PV=nRT$, $PV^x=$ constant becomes $TV^{x-1}=$ constant, so
$T_2=T_1(V_1/V_2)^{x-1}=300\times2^{0.3}\approx\SI{369.3}{K}$. The gas warms by $\Delta T=\SI{69.3}{K}$.

\stepn{2} \textbf{Work.} $W=\dfrac{nR(T_1-T_2)}{x-1}=\dfrac{(8.314)(-69.3)}{0.3}\approx\SI{-1922}{J}$.
Negative, as a compression must be: \SI{1922}{J} of work is done \emph{on} the gas.

\stepn{3} \textbf{Internal energy.} $\Delta U=nC_v\,\Delta T=\tfrac52(8.314)(69.3)\approx\SI{+1441}{J}$,
positive, consistent with the warming.

\stepn{4} \textbf{Heat.} $Q=\Delta U+W=1441+(-1922)=\SI{-480}{J}$. \textbf{Heat leaves the gas}, even though
its temperature rose.

\stepn{5} \textbf{Heat capacity.} $C_x=C_v+\dfrac{R}{1-x}=\tfrac52(8.314)+\dfrac{8.314}{1-1.3}
=20.79-27.71=\SI{-6.93}{J.mol^{-1}.K^{-1}}$. Check: $nC_x\,\Delta T=(1)(-6.93)(69.3)\approx\SI{-480}{J}=Q$.

\stepn{6} \textbf{Read the result.} This is the regime $1<x<\gamma$ of the caution above. The compression
puts in \SI{1922}{J} of work, more than enough to raise $U$ by \SI{1441}{J}, and the remaining
\SI{480}{J} leaves as heat. A negative heat capacity is what happens when work, not heat, drives the
temperature.
\ans{$T_2\approx\SI{369}{K}$, $W\approx\SI{-1922}{J}$, $\Delta U\approx\SI{1441}{J}$, $Q\approx\SI{-480}{J}$, $C_x\approx\SI{-6.93}{J.mol^{-1}.K^{-1}}$}
\end{exbox}

\subsection{Cyclic processes}\label{sec:cyclic}

Every engine works in a cycle: its working substance is taken through a sequence of processes that
returns it to its starting state, ready to go round again. A cycle is a closed loop on the $P$--$V$
diagram, and the first law says something simple about it.

\begin{keybox}{A cycle: \texorpdfstring{$\Delta U=0$}{Delta U = 0}, net work = enclosed area}\label{key:cycle}
Around any closed loop the system returns to its initial state, so every state variable returns to its
initial value: $\Delta U=0$, and the first law gives
\[ Q_{\text{net}}=W_{\text{net}} . \]
The net work is the area under the outgoing part of the loop minus the area under the returning part,
that is, the \textbf{area enclosed by the loop}. It is positive when the loop is traversed clockwise (the
gas expands at higher pressure than it is compressed: an engine) and negative when traversed
anticlockwise (a refrigerator).
\end{keybox}

\begin{figure}[htbp]
\centering
\tikzfig{thermo_rectcycle}
\figcaption{fig:thermo-rectcycle}{A rectangular cycle, $A\to B\to C\to D\to A$, traversed clockwise. The
expansion $B\to C$ happens at the higher pressure and the compression $D\to A$ at the lower, so the gas
does more work than is done on it; the difference is the shaded area. Two legs are at constant volume
(orange, the colour of heat, since no work is done on these legs) and two at constant pressure (green).}
\end{figure}

\begin{exbox}{A rectangular cycle, leg by leg}\label{ex:thermo-rectcycle}
A monatomic ideal gas ($C_v=\tfrac32R$, $C_p=\tfrac52R$) is taken round the cycle of
Figure~\ref{fig:thermo-rectcycle}, between pressures \SI{1e5}{Pa} and \SI{2e5}{Pa} and volumes
\SI{1e-3}{m^3} and \SI{2e-3}{m^3}. Find the work and heat on each leg, and verify the cycle rule.

\stepn{1} \textbf{A useful shortcut.} For an ideal gas $nC_v\,\Delta T=(C_v/R)\,\Delta(PV)$ and
$nC_p\,\Delta T=(C_p/R)\,\Delta(PV)$, since $nR\,\Delta T=\Delta(PV)$. So the heats follow from the change in
$PV$ alone, without knowing $n$ or $T$.

\stepn{2} \textbf{$A\to B$ (constant volume, pressure up).} $W=0$;
$Q=\tfrac32\,V\Delta P=\tfrac32(\num{1e-3})(\num{1e5})=\SI{+150}{J}$.

\stepn{3} \textbf{$B\to C$ (constant pressure, expansion).} $W=P\,\Delta V=(\num{2e5})(\num{1e-3})=\SI{+200}{J}$;
$Q=\tfrac52\,P\Delta V=\SI{+500}{J}$.

\stepn{4} \textbf{$C\to D$ (constant volume, pressure down).} $W=0$;
$Q=\tfrac32(\num{2e-3})(\num{-1e5})=\SI{-300}{J}$.

\stepn{5} \textbf{$D\to A$ (constant pressure, compression).} $W=(\num{1e5})(\num{-1e-3})=\SI{-100}{J}$;
$Q=\tfrac52(\num{-100})=\SI{-250}{J}$.

\stepn{6} \textbf{Totals.} $W_{\text{net}}=200-100=\SI{100}{J}$, which is the enclosed area
$\Delta P\,\Delta V=(\num{1e5})(\num{1e-3})$; $Q_{\text{net}}=150+500-300-250=\SI{100}{J}=W_{\text{net}}$, as
$\Delta U=0$ requires. Note that the gas \emph{absorbed} \SI{650}{J} on the first two legs and
\emph{rejected} \SI{550}{J} on the last two: only the difference became work. Section~\ref{sec:secondlaw}
shows that this rejection is not a design flaw but a law.
\ans{$W_{\text{net}}=Q_{\text{net}}=\SI{100}{J}$; heat absorbed \SI{650}{J}}
\end{exbox}

\begin{trapbox}{``returns to the start'' never means ``nothing happened''}
Around a cycle only the state variables return to their initial values. The transfers do not cancel:
heat was taken in and given out, and net work was done. Writing $W_{\text{net}}=0$ for a cycle confuses
$W$ with $\Delta U$.
\end{trapbox}

% ================================================================
\section{Heat Engines, Refrigerators and the Second Law}\label{sec:secondlaw}

A car engine wastes most of its fuel's energy as heat out of the exhaust and radiator; a coal or nuclear
power station throws away roughly two-thirds of the energy released in its fuel as warm water and
steam. The rectangular cycle above absorbed \SI{650}{J} and rejected \SI{550}{J}. Is this rejection a
failure of engineering, or is it compulsory? The first law cannot tell us: it would be equally satisfied
by an engine that turned all \SI{650}{J} into work. This section shows that the rejection is compulsory,
finds the smallest possible amount, and in doing so meets a second law of nature as fundamental as the
first.

\subsection{What the first law cannot forbid}

The first law is bookkeeping. It forbids creating energy, but it would happily allow a ship that sails by
extracting heat from the sea and turning it into work, or a cup of tea that heats itself up by cooling the
room around it. Energy would balance in both. Neither ever happens. Heat flows from hot to cold of its own
accord, and never the other way; a moving block brought to rest by friction warms the floor, but a warm
floor never cools itself to set a block moving. Processes have a \emph{direction}, and a second,
independent law is needed to say which.

The law was first stated, not for molecules or entropy, but for engines, and it is clearest there. So we
first describe engines and refrigerators, then state the law in two forms and show they say the same
thing.

\subsection{Heat engines and their efficiency}\label{sec:heatengine}

\begin{defbox}{Reservoir, heat engine, efficiency}
A \textbf{heat reservoir} is a body so large that it can give out or take in any amount of heat without
its temperature changing: the atmosphere, the sea, a furnace kept burning. A \textbf{heat engine} is a
device that, working in a cycle, takes heat from a hot reservoir at temperature $T_h$, turns part of it
into work, and gives the rest to a cold reservoir at $T_c$ (Figure~\ref{fig:thermo-engine}). Per cycle let
$Q_h$ be the heat taken in, $Q_c$ the heat rejected and $W$ the work delivered, all three taken
\emph{positive} as magnitudes, with the arrows showing the directions. Its \textbf{efficiency} is the
fraction of the heat taken in that becomes work:
\[
\eta=\frac{W}{Q_h}.
\]
\end{defbox}

For engines and refrigerators it is clearer to work with magnitudes and draw the directions than to carry
the signs of Result~\ref{key:thermo-signs}; the two descriptions agree, since the working substance takes in
$Q=Q_h-Q_c$ net per cycle and does $W$ of work. Petrol and diesel engines (where the ``hot reservoir'' is
burning fuel and the ``cold reservoir'' the atmosphere), steam turbines and jet engines are all heat
engines.

\begin{figure}[htbp]
\centering
\tikzfig{thermo_engine}
\figcaption{fig:thermo-engine}{Every heat engine, reduced to its energy flows per cycle. It takes $Q_h$
from the hot reservoir, delivers work $W$, and gives the remainder $Q_c$ to the cold reservoir. The
first law fixes $W=Q_h-Q_c$; the second law will say that the $Q_c$ arrow can never be removed.}
\end{figure}

\begin{keybox}{Efficiency from the first law}\label{key:efficiency}
The working substance returns to its initial state every cycle, so $\Delta U=0$ and $W=Q_h-Q_c$. Hence
\[
\eta=\frac{W}{Q_h}=\frac{Q_h-Q_c}{Q_h}=1-\frac{Q_c}{Q_h}.
\]
The first law alone allows any $\eta$ up to $1$, which would need $Q_c=0$.
\end{keybox}

For the rectangular cycle of Example~\ref{ex:thermo-rectcycle}, $Q_h=\SI{650}{J}$ (the heat absorbed, on
legs $A\to B$ and $B\to C$) and $W=\SI{100}{J}$, so $\eta=100/650\approx15$ per cent. Note that $Q_h$ is the
heat \emph{absorbed}, summed over the legs where heat enters, and not the net heat.

\subsection{Refrigerators and heat pumps}\label{sec:refrigerator}

Run an engine backwards and it becomes a refrigerator: work is put in, and heat is moved from the cold
reservoir to the hot one (Figure~\ref{fig:thermo-fridge}). In a kitchen refrigerator a compressor does
the work, the cold reservoir is the food compartment and the hot reservoir is the room, which receives the
heat through the coils at the back.

\begin{figure}[htbp]
\centering
\tikzfig{thermo_fridge}
\figcaption{fig:thermo-fridge}{A refrigerator is the engine of Figure~\ref{fig:thermo-engine} with every
arrow reversed: work $W$ is spent to move heat $Q_c$ out of the cold reservoir, and the hot reservoir
receives $Q_h=Q_c+W$. The Clausius statement says the $W$ arrow can never be removed.}
\end{figure}

\begin{defbox}{Coefficient of performance}
A refrigerator is judged by what you want, the heat $Q_c$ removed from the cold space, per unit of what
you pay for, the work $W$:
\[
\text{COP}=\frac{Q_c}{W}=\frac{Q_c}{Q_h-Q_c}.
\]
A \textbf{heat pump} is the same machine used to warm the hot side (a house in winter, drawing heat from the
cold air outside), so its figure of merit is $\text{COP}_{\text{hp}}=Q_h/W=1+\text{COP}$. Unlike an
efficiency, a COP can be, and usually is, greater than $1$; it measures heat \emph{moved}, not energy
created.
\end{defbox}

\subsection{The second law: two statements}

Experience with engines and refrigerators, codified by Kelvin, Planck and Clausius, is summed up in two
statements of what no device can do.

\begin{defbox}{Second law of thermodynamics}\label{def:secondlaw}
\begin{itemize}
\item \textbf{Kelvin--Planck statement (engines).} No process is possible whose \emph{sole} result is the
      absorption of heat from a single reservoir and its complete conversion into work. Every engine must
      reject some heat to a colder body: $Q_c>0$, so $\eta<1$.
\item \textbf{Clausius statement (refrigerators).} No process is possible whose \emph{sole} result is the
      transfer of heat from a colder body to a hotter one. Every refrigerator needs work: $W>0$.
\end{itemize}
\end{defbox}

The words ``sole result'' matter. An ideal gas expanding isothermally turns all the heat it absorbs into
work ($Q=W$, Derivation~\ref{der:isowork}), but that is not the sole result: the gas ends up at a larger
volume and lower pressure, and cannot go on expanding for ever. To repeat the process the gas must be
brought back, and the Kelvin--Planck statement says that a cycle which does this must reject heat. Like
the first law, the second law is a summary of experience, not a theorem derived from mechanics; no
violation has ever been observed. A machine that would break it, such as the ship running on the heat of
the sea, is called a \emph{perpetual-motion machine of the second kind}.

\subsubsection{Why the two statements are equivalent}

The two statements look different, one about engines and one about refrigerators, yet each implies the
other. The proof is short and worth seeing, because the same style of argument proves Carnot's theorem
below: assume a forbidden device exists, connect it to an ordinary one, and show that the combination is
forbidden by the other statement (Figure~\ref{fig:thermo-equivalence}).

\begin{figure}[htbp]
\centering
\tikzfig{thermo_equivalence}
\figcaption{fig:thermo-equivalence}{Each statement of the second law implies the other. The dashed
devices are the hypothetical violators; the grey dashed outline encloses the composite machine. In (a),
the composite takes $Q_h-Q_c$ from the hot reservoir and turns it all into work, leaving the cold
reservoir untouched. In (b), the composite moves $Q_c$ from cold to hot and does nothing else.}
\end{figure}

\begin{derivbox}{Kelvin--Planck and Clausius stand or fall together}\label{der:equivalence}
\stepn{1} \textbf{If Clausius fails, so does Kelvin--Planck.} Suppose a device $X$ could move heat $Q_c$
from the cold reservoir to the hot one with no work (Figure~\ref{fig:thermo-equivalence}a). Run an
ordinary engine between the same reservoirs, adjusted so that it rejects exactly $Q_c$ per cycle while
taking $Q_h$ and delivering $W=Q_h-Q_c$. Per cycle the cold reservoir receives $Q_c$ from the engine and
loses $Q_c$ to $X$: no net change. The hot reservoir loses $Q_h-Q_c$, and $W=Q_h-Q_c$ of work appears.
The composite turns heat from a single reservoir entirely into work with no other result, which the
Kelvin--Planck statement forbids.

\stepn{2} \textbf{If Kelvin--Planck fails, so does Clausius.} Suppose an engine $Y$ could turn heat $Q$ from
the hot reservoir entirely into work $W=Q$ (Figure~\ref{fig:thermo-equivalence}b). Use that work to drive
an ordinary refrigerator, which takes $Q_c$ from the cold reservoir and delivers $Q_c+W=Q_c+Q$ to the hot
one. The hot reservoir loses $Q$ to $Y$ and gains $Q_c+Q$: a net gain of $Q_c$, exactly what the cold
reservoir lost, and no work is left over. The composite moves heat from cold to hot with no other result,
which the Clausius statement forbids.

\stepn{3} So a violation of either statement would allow a violation of the other: logically they are one
law. Use whichever form fits the problem.
\end{derivbox}

\subsection{Reversible processes}\label{sec:reversible}

To find the best possible engine we need the idea of a process that wastes nothing.

\begin{defbox}{Reversible process}
A process is \textbf{reversible} if it can be run backwards so that \emph{both} the system \emph{and} its
surroundings are returned exactly to their initial states. This needs two things: the process must be
quasi-static (Section~\ref{sec:quasistatic}), and it must be free of every dissipative effect, such as
friction, viscosity, electrical resistance, and heat flow across a finite temperature difference. All
other processes are \textbf{irreversible}.
\end{defbox}

Quasi-static is necessary but not enough: a piston compressing a gas slowly against friction is
quasi-static, yet the work done against friction cannot be recovered by running it backwards. Heat flow
across a finite temperature gap is irreversible because, by the Clausius statement, the heat cannot be
made to flow back without outside work. Every real process is irreversible to some degree; reversible
processes are the idealised limit, approached by making changes slowly and frictionlessly and letting
heat flow only across tiny temperature differences.

\subsection{The Carnot cycle}\label{sec:carnot}

In 1824 Sadi Carnot, a French engineer, asked what the most efficient possible engine between two
reservoirs would look like. Its heat must be exchanged reversibly, so heat may flow only while the
working substance sits at the temperature of the reservoir it touches; that means isothermal legs at
$T_h$ and $T_c$. To get from one temperature to the other without exchanging heat at any intermediate
temperature, the connecting legs must exchange no heat at all; that means adiabatic legs.

\begin{figure}[htbp]
\centering
\tikzfig{thermo_carnot}
\figcaption{fig:thermo-carnot}{The Carnot cycle, two isotherms joined by two adiabats and traversed
clockwise, drawn to scale for Example~\ref{ex:thermo-carnot} (monatomic gas, $T_h=\SI{500}{K}$,
$T_c=\SI{300}{K}$). Heat enters only on leg~1 and leaves only on leg~3; the enclosed area is the net work.
For a monatomic gas the adiabats are steep, so the loop is slender.}
\end{figure}

\inwords{The cycle is not one clever design among many; for an engine that exchanges heat with only two
reservoirs it is the \emph{only} reversible one. Any other loop between the same two temperatures would
have to take in or give out heat at some temperature between $T_h$ and $T_c$, and that heat would cross a
finite temperature gap to or from a reservoir.}

\subsubsection{The efficiency of the Carnot engine, derived}

Nothing needs to be taken on trust: with an ideal gas as working substance, every leg of
Figure~\ref{fig:thermo-carnot} can be computed with the results of Section~\ref{sec:processes}.

\begin{derivbox}{Carnot cycle, part 1: heat and work on each leg}\label{der:carnot1}
Let the corners be $A\to B\to C\to D\to A$ as in Figure~\ref{fig:thermo-carnot}, with the isotherms at
$T_h$ (hot) and $T_c$ (cold), and take $n$ moles of ideal gas.

\stepn{1} \textbf{Leg 1, $A\to B$: isothermal expansion at $T_h$.} By Derivation~\ref{der:isowork}, $\Delta U=0$
and
\[
Q_h=W_{AB}=nRT_h\ln\frac{V_B}{V_A}.
\]
This is the only leg on which heat is absorbed.

\stepn{2} \textbf{Leg 2, $B\to C$: adiabatic expansion.} $Q=0$, and Result~\ref{key:adiawork} gives
\[
W_{BC}=\frac{nR\,(T_h-T_c)}{\gamma-1}>0,
\]
as the gas cools from $T_h$ to $T_c$.

\stepn{3} \textbf{Leg 3, $C\to D$: isothermal compression at $T_c$.} Work is done on the gas, $W_{CD}<0$, and
the gas rejects heat of magnitude
\[
Q_c=|W_{CD}|=nRT_c\ln\frac{V_C}{V_D}.
\]

\stepn{4} \textbf{Leg 4, $D\to A$: adiabatic compression.} $Q=0$ and
\[
W_{DA}=\frac{nR\,(T_c-T_h)}{\gamma-1}=-W_{BC}.
\]
\textbf{The two adiabatic works cancel exactly}: one leg pays back what the other delivered. The net work
is decided by the isotherms alone, $W=Q_h-Q_c$, which is also the first law for the whole cycle.
\end{derivbox}

\begin{derivbox}{Carnot cycle, part 2: the volume ratios match}\label{der:carnot2}
The efficiency $\eta=1-Q_c/Q_h$ still contains two logarithms:
\[
\frac{Q_c}{Q_h}=\frac{T_c\ln(V_C/V_D)}{T_h\ln(V_B/V_A)} .
\]
The adiabats make them equal.

\stepn{1} Apply $TV^{\gamma-1}=$ constant to leg $B\to C$: $\;T_hV_B^{\gamma-1}=T_cV_C^{\gamma-1}$.

\stepn{2} Apply it to leg $D\to A$: $\;T_cV_D^{\gamma-1}=T_hV_A^{\gamma-1}$.

\stepn{3} Divide the first equation by the second; the temperatures cancel:
\[
\frac{V_B^{\gamma-1}}{V_A^{\gamma-1}}=\frac{V_C^{\gamma-1}}{V_D^{\gamma-1}}
\quad\Longrightarrow\quad
\frac{V_B}{V_A}=\frac{V_C}{V_D} .
\]
The expansion ratio on the hot isotherm equals the compression ratio on the cold one.

\stepn{4} The two logarithms are therefore equal and cancel:
\[
\boxed{\;\frac{Q_c}{Q_h}=\frac{T_c}{T_h}\;}
\qquad\Longrightarrow\qquad
\boxed{\;\eta_{\text{Carnot}}=1-\frac{T_c}{T_h}\;}
\quad (T\ \text{in kelvin}).
\]

\stepn{5} \textbf{Pause on what happened.} Every trace of the working gas, $n$, $\gamma$ and the volumes,
has vanished. The efficiency depends on the two reservoir temperatures and nothing else. That strongly
suggests the result is universal, and Carnot's theorem below proves it.
\end{derivbox}

\begin{exbox}{The whole Carnot cycle, audited leg by leg}\label{ex:thermo-carnot}
One mole of a monatomic ideal gas ($\gamma=\tfrac53$) runs a Carnot cycle between $T_h=\SI{500}{K}$ and
$T_c=\SI{300}{K}$, doubling its volume on the hot isotherm ($V_A=\SI{1.00}{L}$, $V_B=\SI{2.00}{L}$). Find the
remaining corners, the heats, the works and the efficiency.

\stepn{1} \textbf{Adiabatic corners.} $TV^{2/3}=$ constant gives $V_C/V_B=(T_h/T_c)^{3/2}=(5/3)^{1.5}=2.152$,
so $V_C=\SI{4.30}{L}$; likewise $V_D=2.152\,V_A=\SI{2.15}{L}$. Check: $V_C/V_D=2=V_B/V_A$.

\stepn{2} \textbf{Heats.} $Q_h=RT_h\ln2=(8.314)(500)(0.6931)=\SI{2881}{J}$; $\;Q_c=RT_c\ln2=\SI{1729}{J}$.

\stepn{3} \textbf{Ratio test.} $Q_c/Q_h=1729/2881=0.600=T_c/T_h=300/500$.

\stepn{4} \textbf{Adiabatic legs.} $W_{BC}=R(200)/(2/3)=\SI{2494}{J}$ and $W_{DA}=\SI{-2494}{J}$, cancelling.

\stepn{5} \textbf{Net work and efficiency.} $W=Q_h-Q_c=\SI{1153}{J}$; $\;\eta=1153/2881=0.400=1-300/500$.
\ans{$\eta=40$ per cent, $W=\SI{1153}{J}$ per cycle}
\end{exbox}

\begin{trapbox}{Carnot efficiency needs kelvins, and is a ceiling}
Between \SI{100}{\celsius} and \SI{0}{\celsius} the Carnot efficiency is \emph{not}
$1-0/100=100$ per cent. Convert first: $\eta=1-273/373=26.8$ per cent. And this is a \emph{ceiling}: a
real engine between these temperatures, with its friction and its heat flows across finite temperature
gaps, does strictly worse. Nor can the ceiling be reached by lowering $T_c$ to zero: the coldest
available reservoir is usually the atmosphere or a river, near \SI{300}{K}.
\end{trapbox}

\subsection{Carnot's theorem: no engine can do better}\label{sec:carnot-theorem}

The ideal-gas calculation gives the efficiency of one particular reversible engine. Carnot's theorem
extends it to every engine, using the argument of Derivation~\ref{der:equivalence}.

\begin{derivbox}{A better-than-Carnot engine would violate the second law}\label{der:carnot-theorem}
\stepn{1} Suppose, for contradiction, that an engine $X$ working between the same two reservoirs had
$\eta_X>\eta_{\text{Carnot}}$. Then from the same heat input $Q_h$ it would deliver more work,
$W_X>W_C$, where $W_C$ is the work of a Carnot engine taking in $Q_h$.

\stepn{2} Because the Carnot engine is reversible, it can be run backwards as a refrigerator: fed work
$W_C$, it takes $Q_h-W_C$ from the cold reservoir and delivers $Q_h$ to the hot one, every arrow of
Figure~\ref{fig:thermo-engine} reversed.

\stepn{3} Couple them: let $X$ drive the reversed Carnot engine, using $W_C$ of its output and keeping the
surplus $W_X-W_C>0$.

\stepn{4} Audit the composite over one cycle. The hot reservoir gets back exactly the $Q_h$ that $X$ took,
so it is unchanged. The cold reservoir receives $Q_h-W_X$ from $X$ and loses $Q_h-W_C$ to the
refrigerator: a net loss of $W_X-W_C$, exactly the surplus work.

\stepn{5} So the composite turns heat from a single reservoir entirely into work, which the Kelvin--Planck
statement forbids. The assumption was false:
\[
\boxed{\;\eta\le\eta_{\text{Carnot}}=1-\frac{T_c}{T_h}\ \text{for every engine between the same two reservoirs.}\;}
\]

\stepn{6} Now let $X$ be any \emph{reversible} engine. Then $\eta_X\le\eta_{\text{Carnot}}$, and running the
same argument with the roles swapped ($X$ reversed, driven by the Carnot engine) gives
$\eta_{\text{Carnot}}\le\eta_X$. So every reversible engine between two given reservoirs has
\emph{the same} efficiency, $1-T_c/T_h$, whatever its working substance; the ideal gas was only scaffolding.
An irreversible engine has strictly lower efficiency.
\end{derivbox}

\begin{remarkbox}{Temperature without a thermometer substance}
Because $Q_c/Q_h=T_c/T_h$ holds for every reversible engine, whatever it is made of, Kelvin realised
that it could \emph{define} temperature: the ratio of two temperatures is the ratio of the heats a
reversible engine exchanges with reservoirs at those temperatures. Fixing the triple point of water at
\SI{273.16}{K} then fixes the scale. This \emph{thermodynamic temperature} depends on no property of any
particular substance, and our derivation shows that it coincides with the ideal-gas temperature of
\chref{ch:thermal}, since we used ideal-gas temperatures and got the same ratio.
\end{remarkbox}

\subsubsection{The Carnot refrigerator}

\begin{derivbox}{Coefficient of performance of a Carnot refrigerator}\label{der:carnot-cop}
\stepn{1} Run backwards, the Carnot cycle still has $Q_c/Q_h=T_c/T_h$, since the same legs are traversed in
reverse with the same heats. Divide the numerator and denominator of the COP by $Q_c$:
\[
\text{COP}=\frac{Q_c}{Q_h-Q_c}=\frac{1}{Q_h/Q_c-1}=\frac{1}{T_h/T_c-1}
\quad\Longrightarrow\quad
\boxed{\;\text{COP}_{\text{Carnot}}=\frac{T_c}{T_h-T_c}\;}
\]
and for a heat pump $\text{COP}_{\text{hp}}=T_h/(T_h-T_c)$.

\stepn{2} By the same coupling argument as Carnot's theorem, no refrigerator between the same reservoirs
can have a larger COP.

\stepn{3} \textbf{Read it.} The COP is large when the temperature difference is small and collapses as
$T_h-T_c$ grows. Moving heat ``uphill'' is cheap when the hill is low, which is why a freezer costs more to
run than a refrigerator, and why heat pumps work best in mild winters.
\end{derivbox}

\begin{exbox}{A kitchen refrigerator}\label{ex:thermo-fridge}
A refrigerator keeps its interior at \SI{260}{K} in a room at \SI{300}{K}, and must remove \SI{130}{J} of
heat per second that leaks in through its walls. Find the least electrical power its compressor needs, and
the heat delivered to the room per second.

\stepn{1} \textbf{Best possible COP.} $\text{COP}_{\text{Carnot}}=T_c/(T_h-T_c)=260/40=6.5$.

\stepn{2} \textbf{Least power.} $W=Q_c/\text{COP}=130/6.5=\SI{20}{J}$ per second, that is, \SI{20}{W}.

\stepn{3} \textbf{Heat to the room.} $Q_h=Q_c+W=\SI{150}{J}$ per second. The room gains more heat than the
refrigerator removes from the food; leaving its door open would warm the kitchen, not cool it.
\ans{at least \SI{20}{W}; \SI{150}{W} into the room}
\end{exbox}

% ================================================================
\section{Entropy}\label{sec:entropy}

\scopebox{Entropy is not examined in the NCERT syllabus and only occasionally at JEE Advanced. It is
included because the earlier sections keep pointing at it (free expansion is ``irreversible'', heat
``never flows back''), and because it turns the second law from a list of impossible machines into a
single quantitative statement. The derivations here use only the Carnot results above.}

The two statements of the second law are about what cannot be built. A law of that form is hard to apply
to a process that is not an engine, such as a gas rushing into a vacuum or two blocks coming to the same
temperature. We would like a state variable that measures how far a process is from reversible, as $U$
measures energy. The Carnot cycle hands it to us.

\subsection{From the Carnot cycle to a new state variable}

Write the Carnot result $Q_c/Q_h=T_c/T_h$ with the heats given their signs from the gas's point of view:
the gas takes in $+Q_h$ at $T_h$ and $-Q_c$ at $T_c$. Then
\[
\frac{Q_h}{T_h}+\frac{-Q_c}{T_c}=0 :
\]
around a reversible cycle, the heats taken in, each divided by the temperature at which it is taken in,
add to zero. Clausius showed that the same is true of \emph{any} reversible cycle, by approximating it
with many thin Carnot cycles side by side (their shared adiabatic legs are traversed in opposite
directions and cancel). We state that generalisation without the detailed proof, which needs more care
with limits than this book gives:
\[
\oint\frac{\delta Q_{\text{rev}}}{T}=0\quad\text{around every reversible cycle.}
\]
A quantity whose sum around every closed loop is zero has the same total along every path between two
given states (go out by one path and back by the other: the sum is zero, so the two path sums are equal).
That is exactly the property of a state variable, like the $\Delta U$ of the first law.

\begin{defbox}{Entropy}
The \textbf{entropy} $S$ of a system is the state variable (\si{J/K}, a scalar) whose change between two
equilibrium states is
\[
\boxed{\;\Delta S=S_2-S_1=\int_1^2\frac{\delta Q_{\text{rev}}}{T}\;}
\]
evaluated along \emph{any reversible path} joining the two states, where $\delta Q_{\text{rev}}$ is the heat
taken in on a small step of that path and $T$ the temperature at which it is taken in. Only changes in $S$
are defined this way, just as only changes in potential energy were in \chref{ch:energy}. (In \chref{ch:fluids} the letter $S$ meant surface tension; the two meanings never meet in one chapter, and entropy is $S$ in every book.)
\end{defbox}

Because $S$ is a state variable, $\Delta S$ for an \emph{irreversible} process is found by inventing a
reversible path between the same two end states and computing along that. The actual process may have
quite different heat flows, or none at all.

\begin{derivbox}{Entropy change of an ideal gas}\label{der:entropy-gas}
\stepn{1} On a reversible step, the first law gives $\delta Q_{\text{rev}}=dU+P\,dV=nC_v\,dT+\dfrac{nRT}{V}\,dV$.

\stepn{2} Divide by $T$:
$\dfrac{\delta Q_{\text{rev}}}{T}=nC_v\,\dfrac{dT}{T}+nR\,\dfrac{dV}{V}$. Each term is now a function of a single
variable, which is why the integral does not depend on the path.

\stepn{3} Integrate (two logarithms, \chref{ch:toolkit}, Section~\ref{sec:explog}):
\[
\boxed{\;\Delta S=nC_v\ln\frac{T_2}{T_1}+nR\ln\frac{V_2}{V_1}\;}
\]

\stepn{4} \textbf{Checks.} Isothermal: $\Delta S=nR\ln(V_2/V_1)=Q/T$, since $Q=nRT\ln(V_2/V_1)$. Reversible
adiabatic: $Q=0$ at every step, so $\Delta S=0$; setting the formula to zero gives
$C_v\ln(T_2/T_1)=-R\ln(V_2/V_1)$, that is, $TV^{R/C_v}=TV^{\gamma-1}=$ constant, the adiabatic law of
Derivation~\ref{der:adiabat} recovered by a different road. A reversible adiabatic process is
\emph{isentropic}.
\end{derivbox}

\subsection{The entropy form of the second law}

For an irreversible engine, Carnot's theorem gives $\eta<1-T_c/T_h$, so $Q_c/Q_h>T_c/T_h$ and
\[
\frac{Q_h}{T_h}-\frac{Q_c}{T_c}<0 .
\]
Around a cycle containing irreversible steps, the sum of $Q/T$ (heats with signs, temperatures of the
reservoirs) is \emph{negative}. Clausius generalised this, as before, to the \textbf{Clausius inequality}
$\oint\delta Q/T\le0$, with equality only for reversible cycles. We state its consequence for an isolated
system, which exchanges no heat at all, without the general proof.

\begin{keybox}{Second law: the entropy statement}\label{key:entropy-law}
The entropy of an isolated system never decreases:
\[
\boxed{\;\Delta S_{\text{isolated}}\ge0\;}
\]
with equality only for reversible processes. Applied to a system together with its surroundings (which
together are isolated), $\Delta S_{\text{system}}+\Delta S_{\text{surroundings}}\ge0$. The entropy of one
part \emph{can} fall, as the food in a refrigerator shows, provided the rest gains at least as much.
\end{keybox}

The Clausius statement follows at once. Let heat $Q$ pass from a reservoir at $T_h$ to one at $T_c$. The
hot reservoir loses entropy $Q/T_h$, the cold one gains $Q/T_c$, and
\[
\Delta S_{\text{total}}=Q\left(\frac{1}{T_c}-\frac{1}{T_h}\right),
\]
which is positive if $T_h>T_c$ and negative otherwise. Heat flows from hot to cold, never the reverse,
because only that direction raises the total entropy.

\begin{exbox}{Entropy produced by heat flow and by free expansion}\label{ex:thermo-entropy}
(a) \SI{1000}{J} of heat leaks from a furnace wall at \SI{500}{K} into the air at \SI{300}{K}. Find the
total entropy change. (b) One mole of ideal gas in an insulated container expands freely into an equal
evacuated volume. Find its entropy change.

\stepn{1} \textbf{(a)} $\Delta S=-\dfrac{1000}{500}+\dfrac{1000}{300}=-2.000+3.333=\SI{+1.33}{J/K}$. Positive,
as it must be for a spontaneous process. Had the \SI{1000}{J} instead run a Carnot engine, it would have
delivered $1000(1-300/500)=\SI{400}{J}$ of work and produced no entropy; the leak wastes that opportunity.

\stepn{2} \textbf{(b)} Here $Q=0$ and $W=0$, so $\Delta U=0$ and $T$ is unchanged. Tempting, but wrong, is
$\Delta S=Q/T=0$: the expansion is irreversible, so $Q/T$ does not give $\Delta S$. Instead take a reversible
path between the same end states, an isothermal expansion from $V$ to $2V$:
$\Delta S=nR\ln2=(8.314)(0.693)=\SI{+5.76}{J/K}$. The gas never spontaneously returns to half the box,
because that would lower the entropy of an isolated system.
\ans{(a) $+\SI{1.33}{J/K}$; (b) $+\SI{5.76}{J/K}$}
\end{exbox}

\begin{trapbox}{$\Delta S=Q/T$ only on a reversible path}
For an irreversible process the actual heat divided by the temperature is \emph{less} than $\Delta S$ (it
is zero for the free expansion, while $\Delta S=nR\ln2$). Always compute $\Delta S$ along an imagined
reversible path between the actual end states.
\end{trapbox}

\begin{remarkbox}{Entropy and disorder}
Entropy is often described as ``disorder''. The picture is suggestive but loose. The precise statement,
due to Boltzmann, is $S=k_B\ln\Omega$, where $k_B$ is Boltzmann's constant (\chref{ch:kinetic}) and
$\Omega$ is the number of microscopic arrangements of the molecules consistent with the macroscopic state.
We quote it without proof, since deriving it needs statistical mechanics. When a gas doubles its
volume, each molecule has twice as many places to be, so $\Omega$ grows by a factor $2^{N}$ for $N$
molecules and $S$ by $k_BN\ln2=nR\ln2$, exactly the result of Example~\ref{ex:thermo-entropy}. The second law
then reads: isolated systems move towards macroscopic states that can be realised in overwhelmingly more
ways. A spontaneous decrease is not strictly impossible, only so improbable for $10^{23}$ molecules that
it is never seen; this is why the second law is a statistical law, and why it gives time its direction.
\end{remarkbox}

% ================================================================
\section{Summary}\label{sec:summary-thermo}

The first law counts energy, the second law says which way it flows, and an ideal gas lets both be
computed: every process is a path on the $P$--$V$ diagram, work is the area under it, and the
efficiency of every engine is capped by two temperatures.

\begin{itemize}
\item \textbf{Zeroth law:} thermal equilibrium is transitive, so temperature is a consistent label.
\item \textbf{Work by a gas:} $W=\int P\,dV$, the area under the path; it depends on the path. Heat $Q$
      depends on the path too; the internal energy $U$ is a state variable.
\item \textbf{Signs:} $Q>0$ into the system, $W>0$ done by the system. \textbf{First law:} $Q=\Delta U+W$.
\item \textbf{Ideal gas:} $U$ depends on $T$ only, so $\Delta U=nC_v\,\Delta T$ in every process;
      $C_p-C_v=R$; $\gamma=C_p/C_v$, $C_v=R/(\gamma-1)$.
\item \textbf{Polytropic family} $PV^x=$ constant: $W=nR(T_1-T_2)/(x-1)$, $C_x=C_v+R/(1-x)$; isobaric $x=0$,
      isothermal $x=1$ ($W=nRT\ln(V_2/V_1)$, $Q=W$), adiabatic $x=\gamma$ ($Q=0$, $TV^{\gamma-1}=$ constant,
      expansion cools), isochoric $x\to\infty$ ($W=0$). Slope $-xP/V$: the adiabat is $\gamma$ times steeper
      than the isotherm.
\item \textbf{Cycles:} $\Delta U=0$, $W_{\text{net}}=Q_{\text{net}}=$ enclosed area, positive if clockwise.
\item \textbf{Engines:} $\eta=W/Q_h=1-Q_c/Q_h$, with $Q_h$ the heat absorbed. \textbf{Refrigerators:}
      $\text{COP}=Q_c/W$; heat pumps $\text{COP}_{\text{hp}}=Q_h/W$.
\item \textbf{Second law:} no engine turns all its heat into work (Kelvin--Planck); no heat flows from cold to
      hot unaided (Clausius); the two are equivalent.
\item \textbf{Carnot:} $\eta_{\text{Carnot}}=1-T_c/T_h$ (kelvin), the same for every reversible engine and an
      upper bound for all; $\text{COP}_{\text{Carnot}}=T_c/(T_h-T_c)$.
\item \textbf{Entropy:} $\Delta S=\int\delta Q_{\text{rev}}/T$ along a reversible path. For an ideal gas,
      $\Delta S=nC_v\ln(T_2/T_1)+nR\ln(V_2/V_1)$. The entropy of an isolated system never decreases.
\end{itemize}

% ================================================================
\section{Problem Types}\label{sec:types-thermo}

Thirteen types cover the thermodynamics questions examiners set. Almost all of them turn on two facts:
$\Delta U$ depends only on the end states while $Q$ and $W$ depend on the path, and the process named in
the question (``slowly'', ``suddenly'', ``at constant pressure'') selects the formula.

\archq{First law with the sign convention}
{``Heat $Q$ is supplied and the system expands (or is compressed) by a given amount; find $\Delta U$,'' often
with a phase change.}
{Write $Q=\Delta U+W$ with $Q>0$ for heat in and $W>0$ for work done \emph{by} the system. At constant
pressure, $W=P\,\Delta V$.}
{Adding $W$ instead of subtracting it, which is the first law of the other sign convention applied to this
book's $W$.}
{\SI{1.00}{g} of water at \SI{100}{\celsius} becomes \SI{1671}{cm^3} of steam at a pressure of
\SI{1.013e5}{Pa}; the latent heat is \SI{2256}{J/g}. The increase in internal energy is about:}
{\SI{2256}{J}}
{\SI{2425}{J}}
{\SI{169}{J}}
{\SI{2087}{J}\archans{D}{$W=P\,\Delta V=(\num{1.013e5})(\num{1670e-6})=\SI{169}{J}$ is done by the steam on
the atmosphere, so $\Delta U=Q-W=2256-169=\SI{2087}{J}$. Option (A) ignores the work, (B) adds it (the sign
error), and (C) is the work itself.}}

\archq{First law across two paths}
{``Along one path from $A$ to $B$, heat $Q_1$ is supplied and work $W_1$ done; find $Q$ or $W$ for another path
between the same states.''}
{$\Delta U$ depends only on the end states: compute it on the path where both $Q$ and $W$ are known, then use
$Q=\Delta U+W$ on the other path.}
{Treating $Q$ (or $W$) as if it were the same on both paths. Only $\Delta U$ is.}
{A gas goes from state $A$ to state $B$ along a path on which \SI{600}{J} of heat is supplied and the gas does
\SI{250}{J} of work. Along a different path from $A$ to $B$ the gas does only \SI{100}{J} of work. The heat
supplied along the second path is:}
{\SI{350}{J}}
{\SI{600}{J}}
{\SI{700}{J}}
{\SI{450}{J}\archans{D}{$\Delta U=Q_1-W_1=600-250=\SI{350}{J}$ on either path, so on the second
$Q=\Delta U+W_2=350+100=\SI{450}{J}$. Option (A) is $\Delta U$, not $Q$; (B) assumes $Q$ is path-independent;
(C) adds the work to the old heat.}}

\archq{Net work of a cycle: the enclosed area}
{``A gas is taken clockwise (or anticlockwise) around a closed loop on a $P$--$V$ diagram; find the net work or
the net heat per cycle.''}
{$\Delta U=0$ around a cycle, so $Q_{\text{net}}=W_{\text{net}}=$ enclosed area; positive if clockwise,
negative if anticlockwise.}
{Integrating leg by leg when only the net is asked; or writing $W_{\text{net}}=0$ because the gas ``returns to
the start''.}
{A gas is carried clockwise around a rectangular cycle on a $P$--$V$ diagram between pressures \SI{1e5}{Pa} and
\SI{2e5}{Pa} and volumes \SI{1e-3}{m^3} and \SI{2e-3}{m^3}. The net work done by the gas per cycle is:}
{\SI{100}{J}}
{\SI{50}{J}}
{\SI{200}{J}}
{\SI{0}{J}\archans{A}{The enclosed area is $\Delta P\,\Delta V=(\num{1e5})(\num{1e-3})=\SI{100}{J}$, positive
because the loop is clockwise. Option (B) uses the triangle formula, (C) is the work of the high-pressure
expansion alone, and (D) confuses $W_{\text{net}}$ with $\Delta U$.}}

\archq{Free expansion into a vacuum}
{``A gas in an insulated container expands freely into an evacuated chamber; what happens to its temperature
(or internal energy, or entropy)?''}
{No opposing pressure, so $W=0$; insulated, so $Q=0$; hence $\Delta U=0$, and for an ideal gas $T$ is
unchanged.}
{Expecting cooling because ``expansion cools''. That is true of a reversible adiabatic expansion, where the gas
does work, not of a free one.}
{An ideal monatomic gas, thermally insulated, expands freely into an evacuated chamber and doubles its volume.
Its temperature:}
{falls, because the gas expands}
{is unchanged}
{rises, because the gas speeds up as it rushes in}
{falls to $2^{-2/3}$ of its initial value\archans{B}{$W=0$ and $Q=0$, so $\Delta U=0$, and for an ideal gas
$U$ depends only on $T$. Option (D) applies $TV^{\gamma-1}=$ constant, which holds only for a quasi-static
adiabatic process, and a free expansion is neither quasi-static nor reversible; (A) and (C) have no
energy to pay for a change.}}

\archq{Slow or sudden compression}
{``A gas is \emph{suddenly} (or \emph{slowly}) compressed to a fraction of its volume; find the final pressure
or temperature.'' The adverb is the switch.}
{Slowly, with good thermal contact: isothermal, $PV=$ constant. Suddenly, or insulated: adiabatic,
$PV^\gamma=$ constant and $TV^{\gamma-1}=$ constant.}
{Using the isothermal exponent for a sudden change, or the wrong $\gamma$ for the gas.}
{A monatomic ideal gas ($\gamma=\tfrac53$) at pressure $P_0$ is \emph{suddenly} compressed to one-quarter of its
volume. The final pressure is about:}
{$4P_0$}
{$7.0P_0$}
{$10.1P_0$}
{$16P_0$\archans{C}{Sudden means adiabatic: $P=P_0\,4^{5/3}\approx10.1P_0$. Option (A) is the slow, isothermal
answer; (B) uses the diatomic $\gamma=1.4$; (D) squares the volume ratio.}}

\archq{Slopes of isotherm and adiabat}
{``Compare the slopes of the isothermal and adiabatic curves through a point,'' or ``identify which curve on a
$P$--$V$ diagram is which.''}
{From $PV^x=$ constant, $dP/dV=-xP/V$: the isotherm ($x=1$) has slope $-P/V$, the adiabat ($x=\gamma$) has
$-\gamma P/V$.}
{Thinking the two curves look alike because both fall. The adiabat is always the steeper, by the factor
$\gamma$.}
{At the point where an isothermal and an adiabatic curve of the same gas cross on a $P$--$V$ diagram, the ratio
of the adiabatic slope to the isothermal slope is:}
{$1$}
{$1/\gamma$}
{$\gamma-1$}
{$\gamma$\archans{D}{$(dP/dV)_{\text{adiabatic}}=-\gamma P/V$ and $(dP/dV)_{\text{isothermal}}=-P/V$ at the
same $(P,V)$, so the ratio is $\gamma>1$. Option (B) inverts it; (C) confuses the exponent of $V$ in
$TV^{\gamma-1}$ with the slope.}}

\archq{Molar heat capacity of a given process}
{``A gas follows $VT=$ constant, or $P\propto V$, or $PV^2=$ constant; find its molar heat capacity, or the heat
for a temperature rise.''}
{Rewrite the relation as $PV^x=$ constant (use $PV=nRT$ to eliminate $T$), read off $x$, then
$C_x=C_v+R/(1-x)$ and $Q=nC_x\,\Delta T$.}
{Getting $x$ wrong when converting. $VT=$ constant with $T\propto PV$ gives $PV^2=$ constant, so $x=2$, not
$-1$.}
{One mole of a monatomic ideal gas is taken through a process in which $VT=$ constant. The molar heat capacity
of the gas for this process is:}
{$3R/2$}
{$R/2$}
{$5R/2$}
{$2R$\archans{B}{$T\propto PV$, so $VT=$ constant becomes $PV^2=$ constant and $x=2$:
$C=\tfrac32R+R/(1-2)=\tfrac12R$. Option (A) is $C_v$, ignoring the work; (C) is $C_p$; (D) uses $x=-1$, the
index for $P\propto V$.}}

\archq{Work and heat on one leg}
{``A gas expands along a given path; find $W/Q$, or the fraction of heat converted to work.''}
{Compute $W$ (area, or $nR\,\Delta T$ at constant pressure) and $\Delta U=nC_v\,\Delta T$ separately; then
$Q=\Delta U+W$.}
{Taking $W=Q$, which holds only when $\Delta U=0$ (isothermal).}
{One mole of a monatomic ideal gas expands at constant pressure, doubling its volume from an initial temperature
$T_0$. The ratio $W/Q$ of the work done by the gas to the heat absorbed is:}
{$2/3$}
{$3/5$}
{$2/5$}
{$1$\archans{C}{$W=R\,\Delta T=RT_0$ (the temperature doubles) and $Q=C_p\,\Delta T=\tfrac52RT_0$, so
$W/Q=\tfrac25$. Option (A) is $W/\Delta U$, (B) is $\Delta U/Q$, and (D) forgets that part of the heat raises
$U$.}}

\archq{Carnot efficiency and the kelvin trap}
{``An ideal (Carnot) engine works between two temperatures; find its efficiency, work or rejected heat,'' often
``between the boiling and freezing points of water''.}
{$\eta=1-T_c/T_h$ with temperatures in kelvin; then $W=\eta Q_h$ and $Q_c=Q_h-W$.}
{Using Celsius temperatures; dividing by the wrong temperature; forgetting the ``$1-$''.}
{A Carnot engine operates between the boiling point and the freezing point of water. Its efficiency is closest
to:}
{26.8 per cent}
{36.6 per cent}
{73.2 per cent}
{100 per cent\archans{A}{$\eta=1-273/373=0.268$. Option (B) divides the temperature difference by the cold
temperature, $100/273$; (C) is $T_c/T_h$, the ``$1-$'' forgotten; (D) uses Celsius, $1-0/100$.}}

\archq{Carnot engines in series}
{``Two Carnot engines in series, the heat rejected by the first feeding the second, have equal efficiencies
(or equal works); find the intermediate temperature.''}
{For reversible engines $Q/T$ is the same at every reservoir. Equal efficiency: $1-T/T_h=1-T_c/T$, so
$T=\sqrt{T_hT_c}$. Equal work: the temperature drops are equal, $T=\tfrac12(T_h+T_c)$.}
{Mixing up the two conditions: equal efficiency gives the geometric mean, equal work the arithmetic mean.}
{Two Carnot engines work in series between \SI{800}{K} and \SI{200}{K}, the heat rejected by the first being the
input to the second. If the two have equal efficiency, the intermediate temperature is:}
{\SI{300}{K}}
{\SI{500}{K}}
{\SI{600}{K}}
{\SI{400}{K}\archans{D}{$T/800=200/T$ gives $T=\sqrt{800\times200}=\SI{400}{K}$. Option (B), the arithmetic mean,
is the answer for equal \emph{work}; (A) and (C) have no basis in either condition.}}

\archq{Efficiency of a non-Carnot cycle}
{``A gas goes round a rectangular, triangular or other cycle; find its efficiency.''}
{$\eta=W_{\text{net}}/Q_{\text{in}}$: $W_{\text{net}}$ is the enclosed area; $Q_{\text{in}}$ is the sum of the
heats on the legs where heat is \emph{absorbed} ($Q>0$), each from $nC\,\Delta T$ with the right $C$.}
{Putting the net heat (which equals $W_{\text{net}}$, giving $\eta=1$) or all legs in the denominator.}
{A gas is taken round a cycle: heated at constant volume (pressure rises), then expanded at constant pressure,
then returned to the start along a straight line. In $\eta=W_{\text{net}}/Q_{\text{in}}$ the denominator
$Q_{\text{in}}$ should include the heat of:}
{only the constant-volume leg}
{both the constant-volume and the constant-pressure legs}
{only the constant-pressure leg}
{all three legs, added with their signs\archans{B}{Heat is absorbed on the constant-volume heating ($nC_v\,\Delta T>0$)
and on the constant-pressure expansion ($nC_p\,\Delta T>0$), so both belong in $Q_{\text{in}}$. Option (D)
gives the net heat, which equals $W_{\text{net}}$ and would make $\eta=1$; (A) and (C) each miss a leg that
absorbs heat.}}

\archq{Refrigerators and heat pumps}
{``A refrigerator or heat pump works between two temperatures; find its best COP, or the work for a given
heat.''}
{Refrigerator: $\text{COP}=Q_c/W\le T_c/(T_h-T_c)$. Heat pump: $\text{COP}_{\text{hp}}=Q_h/W\le T_h/(T_h-T_c)$.
Kelvin throughout.}
{Using the refrigerator COP for a heat pump (they differ by $1$), or quoting an efficiency.}
{A heat pump keeps a house at \SI{300}{K} when the outside air is at \SI{270}{K}. The greatest possible amount of
heat delivered to the house per joule of work is:}
{\SI{10}{J}}
{\SI{9}{J}}
{\SI{1.11}{J}}
{\SI{0.10}{J}\archans{A}{$\text{COP}_{\text{hp}}=T_h/(T_h-T_c)=300/30=10$. Option (B) is the refrigerator COP
$T_c/(T_h-T_c)$, which counts the heat taken from outside, not the heat delivered; (C) is $T_h/T_c$;
(D) is the Carnot efficiency.}}

\archq{Entropy change of an irreversible process}
{``Find the entropy change in a free expansion, a mixing, or a heat flow between two bodies.''}
{Entropy is a state variable: replace the actual process by a reversible one between the same end states,
then use $\Delta S=\int\delta Q_{\text{rev}}/T$ (for an ideal gas,
$\Delta S=nC_v\ln(T_2/T_1)+nR\ln(V_2/V_1)$).}
{Using the actual heat: $Q=0$ in a free expansion does not make $\Delta S=0$.}
{One mole of an ideal gas in an insulated container expands freely into vacuum, doubling its volume. The
entropy change of the gas is:}
{zero, because no heat flows}
{$-R\ln2$}
{$+R\ln2\approx\SI{+5.76}{J/K}$}
{$+R\approx\SI{+8.31}{J/K}$\archans{C}{$T$ is unchanged, so take a reversible isothermal expansion from $V$ to
$2V$: $\Delta S=R\ln2=\SI{5.76}{J/K}$. Option (A) uses $Q/T$ for an irreversible process; (B) has the wrong
sign, which would let an isolated system lose entropy; (D) drops the logarithm.}}

% ================================================================
\section{Practice Problems}\label{sec:practice-thermo}

The problems below are of standard examination difficulty. Take $R=\SI{8.314}{J.mol^{-1}.K^{-1}}$.
Answers, with brief working, follow the list.

\begin{enumerate}[itemsep=3pt]
\item A gas absorbs \SI{500}{J} of heat and does \SI{200}{J} of work. Find $\Delta U$. In a second process the
      gas is compressed, \SI{300}{J} of work being done on it, while it loses \SI{100}{J} of heat. Find
      $\Delta U$ again.
\item Half a mole of a diatomic ideal gas is heated at constant pressure from \SI{300}{K} to \SI{400}{K}. Find
      $W$, $\Delta U$ and $Q$.
\item How much heat must be supplied to two moles of an ideal gas at \SI{300}{K} to triple its volume
      isothermally?
\item Air ($\gamma=1.4$) at \SI{300}{K} and atmospheric pressure is compressed adiabatically in a pump until its
      pressure is three times atmospheric. Find its temperature.
\item An ideal monatomic gas is heated along a path on which $P\propto V$. Find its molar heat capacity along
      this path.
\item An inventor claims an engine that, per cycle, absorbs \SI{1000}{J} at \SI{400}{K}, rejects \SI{500}{J} at
      \SI{300}{K} and delivers the rest as work. Is the claim possible? Answer using (a) Carnot's theorem and
      (b) entropy.
\item A freezer at \SI{-18}{\celsius} sits in a kitchen at \SI{30}{\celsius} and must remove heat at
      \SI{100}{W}. Find the least power its compressor can draw.
\item Find the entropy change when \SI{1.0}{kg} of ice melts at \SI{0}{\celsius} ($L_f=\SI{3.34e5}{J/kg}$).
\item Can the temperature of a gas rise when no heat is supplied to it? Can it fall while heat is supplied?
      Give an example of each.
\item Three moles of an ideal gas at \SI{400}{K} are compressed isothermally to half their volume. Find
      the work done \emph{by} the gas and the heat exchanged, with signs.
\item A diatomic gas ($\gamma = 1.4$) is compressed adiabatically to one-eighth of its volume. Starting at
      \SI{300}{K}, find the final temperature.
\item An ideal gas is taken \emph{anticlockwise} round a rectangular cycle on a $P$--$V$ diagram between $P = \SI{1e5}{Pa}$ and \SI{3e5}{Pa}, and $V = \SI{1e-3}{m^3}$ and \SI{4e-3}{m^3}. Find the net work done by the gas per cycle and the net heat it absorbs.
\item A Carnot engine operates between \SI{600}{K} and \SI{300}{K}. Find its efficiency. If it absorbs
      \SI{2000}{J} per cycle, how much work does it deliver?
\item A Carnot engine has an efficiency of 40 per cent with its cold reservoir at \SI{300}{K}. By how much
      must the temperature of the hot reservoir be raised to reach 50 per cent, $T_c$ unchanged?
\item \emph{Conceptual.} A room cannot be cooled by leaving the refrigerator door open. Explain in two
      sentences using the first and second laws.
\item Two moles of a monatomic gas are heated at constant volume from \SI{300}{K} to \SI{500}{K}. Find
      $Q$, $\Delta U$ and $W$. Then repeat for constant pressure, and explain why $Q$ differs.
\item One mole of a monatomic gas expands polytropically, $PV^{1.2}=$ constant, to double its volume,
      starting at \SI{350}{K}. Find the final temperature $T_2$, the work done, the molar heat capacity
      $C_x$ along the process, and $Q$. Is heat absorbed or released? Comment on the sign of $C_x$.
\item A petrol engine burns fuel at about \SI{2000}{K} and exhausts at about \SI{800}{K}. What is the
      \emph{highest} efficiency it could possibly have, and why do real engines fall well short?
\item \emph{Cascade.} A Carnot engine between \SI{600}{K} and \SI{400}{K} feeds the heat it rejects to a
      second Carnot engine between \SI{400}{K} and \SI{300}{K}. Show that the two together have the same
      efficiency as a single Carnot engine between \SI{600}{K} and \SI{300}{K}.
\end{enumerate}

\subsection*{Answers}
{\small
\begin{enumerate}[itemsep=2pt]
\item $\Delta U=Q-W=500-200=\SI{300}{J}$. Second process: $Q=\SI{-100}{J}$, $W=\SI{-300}{J}$, so
      $\Delta U=-100-(-300)=\SI{+200}{J}$.
\item $W=nR\,\Delta T=0.5(8.314)(100)\approx\SI{416}{J}$; $\Delta U=nC_v\,\Delta T=0.5(\tfrac52)(8.314)(100)
      \approx\SI{1039}{J}$; $Q=nC_p\,\Delta T\approx\SI{1455}{J}$, and $416+1039=1455$.
\item $\Delta U=0$, so $Q=W=nRT\ln3=2(8.314)(300)(1.0986)\approx\SI{5480}{J}$.
\item $P^{1-\gamma}T^\gamma=$ constant gives $T_2=T_1(P_2/P_1)^{(\gamma-1)/\gamma}=300\times3^{0.2857}
      \approx\SI{411}{K}$.
\item $P\propto V$ is $PV^{-1}=$ constant, so $x=-1$ and $C=C_v+R/(1-(-1))=\tfrac32R+\tfrac12R=2R
      \approx\SI{16.6}{J.mol^{-1}.K^{-1}}$.
\item (a) Claimed $\eta=500/1000=50$ per cent, but $\eta_{\text{Carnot}}=1-300/400=25$ per cent: impossible.
      (b) Per cycle the engine's entropy returns to its start, while the reservoirs change by
      $-1000/400+500/300=-2.50+1.67=\SI{-0.83}{J/K}$, a decrease of total entropy: impossible.
\item $T_c=\SI{255}{K}$, $T_h=\SI{303}{K}$; $\text{COP}\le255/48=5.31$, so $W\ge100/5.31\approx\SI{18.8}{W}$.
\item Melting is reversible at the melting point, so $\Delta S=Q/T=3.34\times10^5/273\approx\SI{1.22e3}{J/K}$.
\item Yes to both. An adiabatic compression ($Q=0$) heats the gas, the work done on it raising $U$. In a
      polytropic expansion with $1<x<\gamma$, for example $PV^{1.2}=$ constant for air, heat flows in while the
      gas cools, because it does more work than the heat supplied.
\item $W = nRT\ln\tfrac{V_2}{V_1} = 3(8.314)(400)\ln\tfrac12 \approx \SI{-6915}{J}$: the work done
      \emph{by} the gas is negative, since work is done \emph{on} it. Isothermal, so $\Delta U = 0$ and
      $Q = W = \SI{-6915}{J}$: the gas \emph{rejects} that heat.
\item $TV^{\gamma-1}$ is constant: $T_2 = 300\times8^{0.4} \approx 300\times2.30 = \SI{689}{K}$.
      Adiabatic compression heats the gas, as in a bicycle pump.
\item The enclosed area is $(\SI{2e5}{Pa})(\SI{3e-3}{m^3}) = \SI{600}{J}$, and the loop runs anticlockwise (with $V$ horizontal), so the gas is compressed at the higher pressure and $W_{\text{net}}=\SI{-600}{J}$ per cycle: net work is done \emph{on} the gas. $\Delta U = 0$ round the loop, so $Q_{\text{net}} = \SI{-600}{J}$: the gas gives out \SI{600}{J} of heat per cycle, as a refrigerator's working substance does.
\item $\eta = 1 - T_c/T_h = 1-300/600 = 50$ per cent; $W = 0.5\times2000 = \SI{1000}{J}$.
\item $0.4 = 1 - 300/T_h$ gives $T_h = \SI{500}{K}$. For 50 per cent, $T_h = \SI{600}{K}$. Raise
      $T_h$ by \SI{100}{K}.
\item The refrigerator extracts $Q_c$ from its interior but exhausts $Q_h=Q_c + W$ into the room from its
      rear coils (first law), so the net effect on the room is heating by $W$, the compressor's work. The
      second law forbids moving heat from cold to hot without such work, so an open refrigerator
      \emph{warms} the room.
\item \textbf{Constant volume:} $W=0$, so $Q = \Delta U = nC_v\Delta T = 2(\tfrac32R)(200)
      \approx \SI{4988}{J}$. \textbf{Constant pressure:} $\Delta U$ is the \emph{same} \SI{4988}{J}, since
      it depends only on $\Delta T$, but now $W = nR\Delta T = 2R(200) \approx \SI{3326}{J}$ is also done,
      so $Q = \Delta U + W \approx \SI{8314}{J}$. The extra heat pays for the work of expansion, which is
      exactly why $C_p > C_v$.
\item $TV^{x-1}$ is constant with $x=1.2$, so $T_2 = 350(1/2)^{0.2} \approx \SI{304.7}{K}$ (the expansion
      cools the gas, $\Delta T\approx\SI{-45.3}{K}$). $W = \dfrac{nR(T_1-T_2)}{x-1} = \dfrac{8.314(45.3)}{0.2}
      \approx \SI{1883}{J}$, done by the gas. $C_x = C_v + \dfrac{R}{1-x} = \tfrac32R - 5R = -3.5R
      \approx \SI{-29.1}{J.mol^{-1}.K^{-1}}$. Then $Q = nC_x\Delta T = (-29.1)(-45.3)\approx\SI{1318}{J}$,
      \emph{absorbed}. The heat capacity is negative because $1<x<\gamma=\tfrac53$: the gas absorbs heat yet
      cools, doing more work than the heat supplied and drawing the balance from $U$ (check:
      $\Delta U=nC_v\Delta T\approx\SI{-565}{J}=Q-W$).
\item Carnot ceiling: $\eta_{\text{Carnot}} = 1 - 800/2000 = 60$ per cent. Real engines fall short because of
      friction, rapid (non-quasi-static) combustion and heat leaking through the cylinder walls, all
      irreversible.
\item First engine: $\eta_1 = 1 - \tfrac{400}{600}=\tfrac13$, so from heat $Q_h$ taken at \SI{600}{K} it
      rejects $Q_m = \tfrac23 Q_h$ at \SI{400}{K}. The second takes that $Q_m$ and rejects
      $Q_c = Q_m\tfrac{300}{400} = \tfrac23 Q_h\cdot\tfrac34 = \tfrac12 Q_h$ at \SI{300}{K}. Overall
      $\eta = 1 - Q_c/Q_h = \tfrac12 = 1 - \tfrac{300}{600}$, the same as one engine across the whole
      drop: the intermediate temperature cancels.
\end{enumerate}
}

% ================================================================
\section{Challenge Problems}\label{sec:challenge-thermo}

These problems extend the methods of the chapter beyond the standard examination question. Each states
which idea it builds on. Attempt each before reading the solution.

\begin{challengebox}{Isothermal against adiabatic for a diatomic gas (builds on the first law)}
One mole of a diatomic ideal gas ($\gamma=1.4$) at \SI{300}{K} doubles its volume, once isothermally and once
adiabatically. Find the final temperature and the work done in each case, and explain physically which
process does more work.
\challengeans{\emph{Isothermal:} $T_f=\SI{300}{K}$ by definition, and all the heat drawn in becomes work,
$W=nRT\ln2=(8.314)(300)(0.693)\approx\SI{1729}{J}$. \emph{Adiabatic:} no heat enters, so the work comes out
of internal energy and the gas cools:
\[
T_f=T_i\left(\frac{V_i}{V_f}\right)^{\gamma-1}=300\,(0.5)^{0.4}\approx\SI{227.4}{K},
\qquad
W=\frac{nR(T_i-T_f)}{\gamma-1}=\frac{(8.314)(72.6)}{0.4}\approx\SI{1510}{J}.
\]
The isothermal expansion does more work because the surroundings keep supplying heat to hold its
temperature, and so its pressure, up; the adiabatic gas has only its own internal energy to spend.}
\end{challengebox}

\begin{challengebox}{How good is the rectangular cycle? (builds on cycles and Carnot's theorem)}
For the rectangular cycle of Example~\ref{ex:thermo-rectcycle} (pressures \SI{1e5}{Pa} and \SI{2e5}{Pa},
volumes \SI{1}{L} and \SI{2}{L}, traversed clockwise): (a) find its efficiency with a monatomic and with a
diatomic working gas; (b) find the ratio of the highest to the lowest temperature reached, and the Carnot
efficiency between those temperatures; (c) explain why the cycle falls so far short.
\challengeans{(a) $W_{\text{net}}=\SI{100}{J}$ in both cases. Heat is absorbed on $A\to B$ ($Q=(C_v/R)V\Delta P$)
and $B\to C$ ($Q=(C_p/R)P\Delta V$). Monatomic: $Q_{\text{in}}=\tfrac32(100)+\tfrac52(200)=\SI{650}{J}$, so
$\eta=100/650\approx15.4$ per cent. Diatomic: $Q_{\text{in}}=\tfrac52(100)+\tfrac72(200)=\SI{950}{J}$, so
$\eta\approx10.5$ per cent; more of the heat goes into the extra (rotational) internal energy.
(b) $T\propto PV$: the hottest state is $C$ ($PV=\SI{400}{J}$) and the coldest $A$ ($PV=\SI{100}{J}$), a ratio of
4. A Carnot engine between them would have $\eta=1-\tfrac14=75$ per cent.
(c) Carnot's bound is reached only if all heat is taken in at the highest temperature and given out at the
lowest. Here most of the heat is absorbed while the gas is far colder than the hottest state $C$ (on $A\to B$ its temperature only rises from
that of $A$ to twice that), and rejected while it is far hotter than the coldest state $A$. In a real engine those heats would cross
large temperature gaps to and from the reservoirs, which is irreversible; the efficiency pays for it.}
\end{challengebox}

\begin{challengebox}{The polytropic surprise (builds on molar heat capacities)}
A diatomic ideal gas is taken along the process $PV^{1.2}=$ constant. Find its molar heat capacity for this
process, and interpret the sign.
\challengeans{With $x=1.2$,
\[
C=C_v+\frac{R}{1-x}=\frac{5R}{2}+\frac{R}{1-1.2}=\frac{5R}{2}-5R=-\frac{5R}{2}\approx\SI{-20.8}{J.mol^{-1}.K^{-1}}.
\]
A \emph{negative} heat capacity: when the gas expands along this path and takes in heat, its temperature
\emph{falls}. Between the isotherm ($x=1$) and the adiabat ($x=\gamma=1.4$), the expanding gas does more work
than the heat it receives, and the deficit comes out of its internal energy. (Stars also have a negative heat
capacity, for a different reason: as a star radiates energy away, gravity pulls it tighter and its core
\emph{heats up}.)}
\end{challengebox}

\begin{challengebox}{Two engines in series (builds on Carnot's theorem)}
A Carnot engine runs between \SI{600}{K} and an intermediate reservoir at $T$, and a second Carnot engine
takes all the heat rejected by the first and runs between $T$ and \SI{300}{K}
(Figure~\ref{fig:thermo-cascade}). The two produce equal work. Find $T$, and compare the combined efficiency
with that of a single Carnot engine between \SI{600}{K} and \SI{300}{K}.
\challengeans{For reversible engines $Q_h/600=Q_m/T=Q_c/300$; call the common value $\Delta S_0$. Then
$W_1=Q_h-Q_m=\Delta S_0(600-T)$ and $W_2=Q_m-Q_c=\Delta S_0(T-300)$. Equal work requires equal temperature drops:
$T=\tfrac12(600+300)=\SI{450}{K}$, the arithmetic mean. The combined efficiency is
\[
\eta=\frac{W_1+W_2}{Q_h}=1-\frac{Q_c}{Q_h}=1-\frac{300}{600}=50\text{ per cent},
\]
the same as a single Carnot engine across the whole range. Staging neither helps nor hurts: the intermediate
temperature cancels, and Carnot's theorem forbids beating $1-T_c/T_h$ however the drop is divided.
(The common value $\Delta S_0=Q/T$ is the entropy that flows through the chain unchanged, Section~\ref{sec:entropy}.)}
\end{challengebox}

\begin{figure}[htbp]
\centering
\tikzfig{thermo_cascade}
\figcaption{fig:thermo-cascade}{Two Carnot engines in series. The heat $Q_m$ rejected by engine~1 into the
intermediate reservoir is all taken up by engine~2, so the intermediate reservoir neither gains nor loses
energy per cycle. Because $Q/T$ is the same at every reservoir of a reversible chain, each engine's work is
proportional to its own temperature drop.}
\end{figure}

\begin{challengebox}{The most work from two bodies (builds on entropy)}
Two identical blocks, each with heat capacity $ms=\SI{1000}{J/K}$ (mass times specific heat, as in \chref{ch:thermal}) independent of temperature, are at
\SI{400}{K} and \SI{100}{K}. (a) If they are simply put in contact, what is the final temperature, and how
much work is obtained? (b) Instead a reversible engine runs between them, taking heat from the hot block and
rejecting it to the cold one, until they reach a common temperature $T_f$. Find $T_f$ and the work delivered.
\challengeans{(a) Energy conservation gives $T_f=\tfrac12(400+100)=\SI{250}{K}$ and no work at all. The total
entropy rises by $ms\ln(250/400)+ms\ln(250/100)=1000(-0.470+0.916)\approx\SI{+446}{J/K}$.
(b) A reversible process leaves the total entropy unchanged. The blocks change by
$ms\ln(T_f/400)+ms\ln(T_f/100)=0$, so $T_f^2=400\times100$ and $T_f=\sqrt{400\times100}=\SI{200}{K}$. The work is
the heat lost by the hot block minus the heat gained by the cold one:
\[
W=ms(400-T_f)-ms(T_f-100)=ms\left(\sqrt{400}-\sqrt{100}\right)^2=1000\times(20-10)^2=\SI{1.0e5}{J}.
\]
The engine leaves the blocks colder (\SI{200}{K} instead of \SI{250}{K}) because it has carried
$2ms(\SI{50}{K})=\SI{1.0e5}{J}$ of their internal energy off as work. Direct contact throws away exactly this
opportunity, and the entropy it creates measures the loss.}
\end{challengebox}
